\documentclass[11pt]{article}
\usepackage[margin=1in]{geometry}
\usepackage{amsmath,amssymb,amsthm,mathtools}
\usepackage{booktabs,array}
\usepackage{microtype}
\usepackage[numbers,sort&compress]{natbib}
\usepackage[colorlinks=true,allcolors=blue]{hyperref}
\newtheorem{theorem}{Theorem}[section]
\newtheorem{proposition}[theorem]{Proposition}
\newtheorem{lemma}[theorem]{Lemma}
\newtheorem{corollary}[theorem]{Corollary}
\theoremstyle{definition}
\newtheorem{definition}[theorem]{Definition}
\newtheorem{assumption}[theorem]{Assumption}
\newtheorem{example}[theorem]{Example}
\theoremstyle{remark}
\newtheorem{remark}[theorem]{Remark}
\newcommand{\R}{\mathbb R}
\newcommand{\Q}{\mathbb Q}
\newcommand{\Z}{\mathbb Z}
\newcommand{\norm}[1]{\lVert#1\rVert}
\DeclareMathOperator{\conv}{conv}
\DeclareMathOperator{\argmin}{argmin}
\title{Encoding and Calibration of Exact Norm Penalties\allowbreak\ in Mixed-Integer Convex Optimization}
\author{}
\date{}
\begin{document}
\maketitle
\begin{abstract}
Exact penalty methods move equality constraints that link the parts of a
mixed-integer model into the objective as a norm of their violation. A
finite coefficient that closes the duality gap is known to exist for
mixed-integer linear programs and, under compactness and a Slater-type
condition on the integer slices, for mixed-integer convex problems. But
existence does not say how large the coefficient must be or how hard it
is to find. We study both questions for coefficients written in ordinary
binary encoding. First, the polynomial bit length known for linear
constraints is lost once convex quadratic constraints are kept in the
subproblem: a family with one binary variable, $d+1$ continuous
variables and $O(d\log(d+1))$ input bits has least penalty $2^{2^d-1}$,
even with an optimal multiplier and although every integer slice has
points at which all inequalities hold with slack $1/8$. Second, this
growth needs both many continuous variables and many constraints that
are nonlinear in them. For quadratic data that are convex in the
continuous variables, affine linking equalities, finite bounds and a
Slater condition on every integer slice that can satisfy the linking
equalities, some integer
coefficient with $(N+1)2^{O(n)}$ bits is exact at multiplier zero for the
$\infty$- and $1$-norms and preserves the set of minimizers, where $N$ is
the input length and $n$ the number of continuous variables. This bound
is tight up to the constant in the exponent, and $N^{O(k+1)}$ bits
suffice when at most $k$ inequalities are nonlinear in the continuous
variables. Third, unless $\mathrm P=\mathrm{NP}$, no polynomial-time
algorithm approximates the least penalty within any polynomial factor,
even on a binary box with one equality, a known unique optimizer and a
known sufficient coefficient. Finally, for finitely many compact convex
slices with continuous convex objectives and affine residual maps, if the
right-hand side is sampled from a rational grid with at
least $4mK/\epsilon$ points per coordinate in a cube of radius $\sigma$,
then with probability at least $1-\epsilon$ either the perturbed problem
is infeasible or the $\infty$-norm coefficient $4mKM/(\sigma\epsilon)$ is
exact for it at multiplier zero. Here $m$ counts the linking equalities,
$K$ bounds the number of integer slices and $M$ is a supplied bound on
the objective range; the order $K/\epsilon$ is necessary.
\end{abstract}
\section{Introduction}\label{sec:intro}
Many mixed-integer models keep most of their constraints in a set $X$
and couple blocks, stages or scenarios through equalities $r(x)=0$.
Examples are linking constraints between subsystems, coupling
constraints in market models, copy constraints between a stage and its
state, and nonanticipativity constraints. We call the constraints that
define $X$ the \emph{native constraints} and the equalities $r(x)=0$
the \emph{relaxed} (or \emph{linking}) \emph{equalities}. The vector $x$ collects the integer and
continuous variables. The problem is
\[
 v=\min\{f(x):x\in X,\ r(x)=0\},
\]
and moving the relaxed equalities into the objective with a multiplier
$\lambda\in\R^m$ and a norm penalty with coefficient $\rho\ge0$ gives
\[
 L_\rho(\lambda)=\min_{x\in X}\{f(x)+\lambda^Tr(x)+\rho\norm{r(x)}\},
 \qquad
 D_\rho=\sup_{\lambda\in\R^m}L_\rho(\lambda).
\]
Section~\ref{sec:geometry} states the setting precisely
(compact $X$, continuous $f$ and $r$, a fixed norm). Weak duality gives
$L_\rho(\lambda)\le D_\rho\le v$. For $\rho=0$, $D_0$ is the ordinary
Lagrangian dual, which for mixed-integer problems can leave a gap
$D_0<v$. The norm term, which gives the sharp augmented Lagrangian (not
the squared term of the classical augmented Lagrangian), can close this
gap with a finite coefficient. Such a coefficient is an
\emph{exact penalty}.

Exact penalty functions were introduced by \citet{Eremin1966} and
\citet{Zangwill1967}, and \citet{Burke1991} unified their theory.
\citet{Rockafellar1974} showed that augmented Lagrangians remove duality
gaps in nonconvex programming. \citet{RockafellarWets1998} and
\citet{HuangYang2003} developed the connection between augmented
Lagrangian duality and exact penalization. For mixed-integer problems,
finite exact penalties are known for bounded pure integer programs,
rational MILPs, convex MIQPs, convex objectives over polyhedral
mixed-integer sets, and mixed-integer nonlinear problems under
regularity conditions on the integer slices
\citep{BolandEberhard2015,FeizollahiAhmedSun2017,GuAhmedDey2020,Bhardwaj2024,LefebvreSchmidt2025};
Section~\ref{sec:related} states these results precisely.

We use the following terms; Definition~\ref{def:exactness} states them
formally. A coefficient $\rho$ is \emph{value-exact} if $D_\rho=v$, that
is, if the penalized dual with the best multiplier has the optimal value.
It is \emph{exact at $\lambda$} if $L_\rho(\lambda)=v$ for that fixed
multiplier. It is \emph{minimizer-exact at $\lambda$} if the minimizers of
$f+\lambda^Tr+\rho\norm r$ over $X$ are exactly the optimal solutions of
the original problem. Minimizer-exactness at $\lambda$ implies exactness
at $\lambda$, which implies value-exactness; Section~\ref{sec:geometry}
shows that the converses fail. The \emph{least penalty}
$\rho_*=\inf\{\rho\ge0:D_\rho=v\}$ is, when finite, the smallest
value-exact coefficient, with the multiplier optimized. The
\emph{zero-multiplier threshold} $\rho_*(0)$ is, when finite, the
smallest coefficient that is exact at $\lambda=0$. We call a coefficient
\emph{sufficient} if it is value-exact, or exact at $\lambda=0$ when the
zero-multiplier threshold is meant.

\paragraph{Why the size of the coefficient matters.}
Exact norm penalties are used in decomposition methods, and each of them
must choose the coefficient. In cutting-plane methods for multistage stochastic MILPs,
augmented Lagrangian cuts with a norm as augmenting function give
Lipschitz lower approximations of cost-to-go functions; they are tight
for a suitable multiplier and large enough $\rho$, but arbitrarily
large $\rho$ may be needed near discontinuities of the value function
\citep{AhmedCabralCosta2022}.
\citet{ZhangSun2022} penalize copy constraints between a state and its
local copy, assume that this penalty reformulation is exact (their
Assumption~2), and their iteration bounds depend on the penalty sizes. For Lipschitz
regularization in stochastic dual dynamic programming, which penalizes
such copy constraints by a norm, a recent survey calls the choice of a
sufficiently large coefficient an open challenge in practice
\citep{FullnerRebennack2025}.
In $\ell_1$ augmented Lagrangian and ADMM decompositions of two-block
MILPs, the penalty that guarantees an $\varepsilon$-accurate solution
grows like $1/\varepsilon$, and such a large penalty slows the method
\citep{SunSunYin2024}; in bundle methods for generalized augmented
Lagrangian duals, adjusting the penalty parameter was the delicate part
\citep{CordovaOliveiraSagastizabal2022}; and penalty alternating
direction heuristics increase weighted $\ell_1$ penalties on copy
constraints in an outer loop \citep{GeisslerMorsiScheweSchmidt2017}.
In $\ell_1$ augmented Lagrangian pricing of nonconvex electricity
markets, the coefficient is part of the price signal, and
\citet{WangHesamzadehZhaoKronqvist2026} seek a valid coefficient as close
to the least one as is computationally feasible. Large coefficients
also cost precision: floating-point MIP solvers suffer from wide ranges
of data magnitudes \citep{KlotzNewman2013}, and quantum annealers rescale
QUBO coefficients to a fixed range, so a large penalty compresses the
gaps between objective values \citep{KrpanPovhPucher2025}.

These uses raise two questions about a sufficient coefficient. We
measure coefficients in ordinary binary encoding: a rational number is
given by the binary representations of its numerator and its positive
denominator.
\begin{itemize}
\item \emph{Representation.} How many bits does a sufficient coefficient
 need, as a function of the encoding length of the instance?
\item \emph{Calibration.} How closely can the least penalty $\rho_*$ be
 approximated, from above, by an efficient algorithm?
\end{itemize}
Bit length and numerical size are different measures: a coefficient with
polynomially many bits may still be exponentially large as a number. A
coefficient whose bit length is superpolynomial cannot even be written
down in polynomial time.

Both questions are posed in the conclusion of
\citet{LefebvreSchmidt2025}. They note that an exact penalty of
polynomial size is known only for MIQPs and ask whether this extends to
classes such as quadratically constrained quadratic problems. They also
observe that the coefficient computed by their polynomial-time method is
``too large to be useful in practical applications'', and they ask
whether the smallest penalty that closes the duality gap can be computed
in polynomial time, conjecturing that it cannot. We answer the first
question negatively for convex quadratic native constraints when both
the continuous dimension and the number of nonlinear constraints grow,
and positively when either is fixed, under explicit variable bounds and
a Slater condition on the integer slices (Assumption~\ref{ass:model});
for nonconvex quadratic constraints the question remains open
(Section~\ref{sec:discussion}). We confirm the conjecture unless
$\mathrm P=\mathrm{NP}$, even for approximation within any polynomial
factor.

\subsection{Main results}
Table~\ref{tab:results} summarizes the results. Throughout, $N$ is the
encoding length of an instance, $n$ the number of continuous variables
and $m$ the number of relaxed equalities.

\begin{table}[!t]
\centering
\small
\renewcommand{\arraystretch}{1.15}
\begin{tabular}{@{}>{\raggedright\arraybackslash}p{0.12\linewidth}
 >{\raggedright\arraybackslash}p{0.35\linewidth}
 >{\raggedright\arraybackslash}p{0.33\linewidth}
 >{\raggedright\arraybackslash}p{0.12\linewidth}@{}}
\toprule
Result & Setting and hypotheses & Conclusion & Reference\\
\midrule
Lower bound &
One binary and $n=d+1$ continuous variables; linear objective;
$d-1$ convex quadratic native inequalities ($d\ge2$); coefficients from a
fixed finite set; sparse encoding length $O(d\log(d+1))$; points with slack $1/8$
on every integer slice &
$\rho_*=2^{2^d-1}$, $\rho_*(0)=2^{2^d}$; every value-exact rational
coefficient has at least $2^d$ numerator bits; a gap at most
$\epsilon<1$ needs $2^d-O_\epsilon(1)$ bits &
Thm.~\ref{thm:lower}, Cor.~\ref{cor:accuracy}\\
\addlinespace
Lower bound in $N$ &
Same chain with seed $a_1\ge2^{-t}$, $t$ of order $N$ &
least penalty with $\Omega(N2^n)$ bits &
Rem.~\ref{rem:seed}\\
\addlinespace
Upper bound in $n$ &
Assumption~\ref{ass:model}: degree-two data, convex in the continuous
variables for each integer assignment; finite bounds; nonempty feasible
set; refined Slater condition; affine relaxed equalities &
an integer $\rho$ with $(N+1)2^{O(n)}$ bits, minimizer-exact at
$\lambda=0$ for both $\norm\cdot_\infty$ and $\norm\cdot_1$ &
Thm.~\ref{thm:generalupper}\\
\addlinespace
Upper bound in $k$ &
Assumption~\ref{ass:model} with at most $k$ native inequalities nonlinear
in the continuous variables; any number of affine rows &
the same with $N^{O(k+1)}$ bits; for fixed $n$ or fixed $k$ such a
coefficient can be printed in polynomial time &
Thm.~\ref{thm:fixedcount}, Cor.~\ref{cor:printing}\\
\addlinespace
Calibration &
Binary box, one equality, linear objective; known unique optimizer;
$\rho=1$ is exact at $\lambda=0$ &
unless $\mathrm P=\mathrm{NP}$, no polynomial-time algorithm finds
$\widehat\rho$ with $\rho_*\le\widehat\rho\le p(N)\rho_*$, and likewise
for $\rho_*(0)$; testing value-exactness of $\rho=1/3$ is
coNP-complete &
Thm.~\ref{thm:polyfactor}, Thm.~\ref{thm:coNP}\\
\addlinespace
Strong hardness &
Stable-set native set; coefficients in $\{-1,0,1\}$ &
$\rho_*=\alpha(G)$, so computing $\rho_*$ is strongly NP-hard &
Prop.~\ref{prop:stableset}\\
\addlinespace
Perturbation &
At most $K$ compact convex slices; continuous convex objectives with
supplied range $M$; affine residuals; right-hand side sampled from a
rational grid with at least $\lceil4mK/\epsilon\rceil$ points per
coordinate in a cube of radius $\sigma$ &
with probability at least $1-\epsilon$, the sampled problem is infeasible
or $4mKM/(\sigma\epsilon)$ is exact at $\lambda=0$ for
$\norm\cdot_\infty$; the order $K/\epsilon$ is necessary &
Thm.~\ref{thm:gridsmoothing}, Sec.~\ref{sec:numerical-size}\\
\bottomrule
\end{tabular}
\caption{Informal summary of the main results. $\rho_*$: least penalty
(multiplier optimized); $\rho_*(0)$: zero-multiplier threshold;
$p$: any fixed polynomial.}
\label{tab:results}
\end{table}

\paragraph{A doubly exponential lower bound.}
Theorem~\ref{thm:lower} gives, for each chain length $d$, an instance with
one binary variable, $n=d+1$ continuous variables, a linear objective,
one relaxed equality, and convex quadratic and affine native constraints
with coefficients from a fixed finite set; its sparse encoding length is
$O(d\log(d+1))$. Its least penalty is $2^{2^d-1}$, even after optimizing
the multiplier, so every exact rational coefficient, in any of the three
senses, needs at least $2^d$ numerator bits. This holds although every
integer slice has points with slack at least $1/8$, and a fixed additive
accuracy below the objective range does not help
(Corollary~\ref{cor:accuracy}).

The mechanism is simple. A chain of convex quadratic inequalities
$a_{i+1}\ge a_i^2$ lets the improving integer assignment come within
$\delta_d=2^{-2^d}$ of satisfying the relaxed equality, but never reach
it. The other assignment has residuals of both signs at zero cost. A
convex combination of its point with residual $-1$ and the improving
point, weighted so that the residuals average to zero (a \emph{balanced
mixture}, Section~\ref{sec:geometry}), cancels every multiplier, and its
average penalized cost is negative unless $\rho\ge1/(2\delta_d)$. The numbers grow fast. For $d=5$, with six
continuous variables, the improving assignment violates the relaxed
equality by $\delta_5=2^{-32}\approx2.3\cdot10^{-10}$, below common
feasibility tolerances, and the least penalty is $2^{31}$. For $d=10$ it
is $2^{1023}$, the largest power of two in IEEE double precision; for
$d=11$ it exceeds the double-precision range. Thus finite exactness with
convex quadratic native constraints does not come with the polynomial
encoding guarantee known for linear native constraints. The obstruction
lies in the representation of the penalty, not in the original problem,
whose optimal value is immediate.

\paragraph{Encoding upper bounds.}
For the positive results we fix the integer variables and look at the
resulting \emph{slices} of the native set. A slice is
\emph{equality-feasible} if it contains a point satisfying the relaxed
equalities. The \emph{refined Slater condition} requires, for every
equality-feasible slice, a point of the slice that satisfies the relaxed
equalities and strictly satisfies every native inequality that is
nonlinear in the continuous variables; inequalities affine in the
continuous variables may hold with equality, and nothing is assumed on
the other slices. Assumption~\ref{ass:model} asks for degree-two data
that are convex in the continuous variables for each integer
assignment, explicit finite bounds on all variables, affine relaxed
equalities, a nonempty feasible set and the refined Slater condition.
Under this assumption, Theorem~\ref{thm:generalupper} gives an integer
coefficient with $(N+1)2^{O(n)}$ bits, and
Theorem~\ref{thm:fixedcount} gives one with $N^{O(k+1)}$ bits when at
most $k$ native inequalities are nonlinear in the continuous variables;
the number of affine rows is arbitrary. In both cases one coefficient
is minimizer-exact at multiplier zero for the infinity norm and the
one-norm. By Remark~\ref{rem:seed}, the worst case of the first bound
is $N2^{\Theta(n)}$ bits.

Both proofs use Lemma~\ref{lem:slices}, which reduces exactness at
multiplier zero to a multiplier bound on each equality-feasible slice
and a lower bound on the residual of each equality-infeasible slice;
effective bounds from real algebraic geometry control both uniformly
over all integer assignments. Consequently, within the class of
Assumption~\ref{ass:model}, coefficients with superpolynomially many
bits require both a growing continuous dimension and a growing number
of nonlinear native inequalities; the lower-bound instance has both.
For fixed $n$, and separately for fixed $k$, a conservative coefficient
can be printed in polynomial time (Corollary~\ref{cor:printing}). These
are existence bounds: the coefficient may be numerically huge, the
constants are not evaluated, and nothing is said about solving the
penalized subproblem.

\paragraph{Hardness of calibration.}
Theorem~\ref{thm:polyfactor} considers pure binary linear programs with
one equality: the native set is a binary box and the objective is
linear. On these instances the unique optimal solution is known and
$\rho=1$ is exact at multiplier zero. Even so, unless
$\mathrm P=\mathrm{NP}$, no polynomial-time algorithm can output a
sufficient coefficient within any polynomial factor of the least
penalty, or of the zero-multiplier threshold; Section~\ref{sec:calibration}
extends this to factors $2^{N^{1-\epsilon}}$ for every fixed
$\epsilon\in(0,1)$. Deciding whether the
coefficient $\rho=1/3$ is value-exact is coNP-complete
(Theorem~\ref{thm:coNP}). The instances encode \textsc{Subset Sum}: when the target sum is
attainable, some integer assignment that lowers the objective has a
residual of magnitude $1$; otherwise every such assignment has a residual
of magnitude at least $\kappa-1$, for a scaling integer $\kappa$ that we
take exponentially large. The hardness is weak, since a pseudo-polynomial
dynamic program computes $\rho_*$; a stable-set variant with
coefficients in $\{-1,0,1\}$ makes exact computation strongly NP-hard
(Proposition~\ref{prop:stableset}). The message is that a safe
coefficient can be easy to find while a nearly least one is hard to
certify.

\paragraph{Right-hand-side perturbations.}
Theorem~\ref{thm:gridsmoothing} considers at most $K$ compact convex
slices with continuous convex objectives of range $M$ and affine
residual maps. If the right-hand side is sampled from a rational grid
with at least $\lceil4mK/\epsilon\rceil$ points per coordinate in a cube
of radius $\sigma$, then with probability at least $1-\epsilon$ the
sampled problem is infeasible or the infinity-norm coefficient
$4mKM/(\sigma\epsilon)$ is exact at multiplier zero. If the number of grid points per coordinate is the smallest power of
two that is at least $\lceil4mK/\epsilon\rceil$, the grid points and the
coefficient have encoding length polynomial in the supplied rational
data, in $m$ and in $\log K$ (Remark~\ref{rem:gridencoding}). The mechanism is that
large penalties arise only when the right-hand side lies near the
boundary of the set of residuals attainable on some slice; a random
right-hand side avoids these boundaries with high probability, and a
convex repair argument turns the distance into a coefficient. The
guarantee concerns the perturbed problem, whose optimal integer
assignment may differ from the original one. The coefficient is
numerically of order $K$, which can be exponential in the number of
integer variables, and Section~\ref{sec:numerical-size} shows that the
order $K/\epsilon$ is necessary and that the expected least penalty can
be infinite.

\subsection{Related work}\label{sec:related}
\paragraph{Exact penalties and augmented Lagrangian duality.}
The dual of \citet{Rockafellar1974} maximizes jointly over the
multipliers and the penalty parameter; we fix $\rho$, optimize
$\lambda$, and ask for the least $\rho$ at which the optimized value is
exact.\footnote{Theorem, section and example numbers of preprints and
author manuscripts refer to the versions identified in the references.}
Section~\ref{sec:geometry} relates the threshold $\rho_*(\lambda)$ to
the exact penalty representations of the sharp Lagrangian in
\citet[Chapter~11]{RockafellarWets1998}. In the study of global
saddle points of augmented Lagrangians for continuous cone-constrained
problems, \citet{Dolgopolik2018} uses the term ``least exact penalty
parameter''; our least penalty $\rho_*$ is the analogous quantity for
mixed-integer problems, taken after optimizing the multiplier.

\paragraph{Exactness and encoding for mixed-integer problems.}
\citet{BolandEberhard2015} prove finite-penalty exactness for bounded
pure integer programs under their assumptions on the augmenting
function, and \citet{FeizollahiAhmedSun2017} prove norm-penalty exactness
for rational MILPs. \citet[Theorem~11]{GuAhmedDey2020} show that for a
convex quadratic objective over a rational mixed-integer polyhedron some
exact norm penalty has polynomial encoding length.
\citet[Theorems~3.1--3.3]{Bhardwaj2024} give value-function proofs of
existence with data-dependent penalty parameters for rational MILPs and
for MIQPs with bounded integer variables. For convex objectives over
polyhedral mixed-integer sets they prove exactness when the integer
variables are bounded on the feasible set or when the objective and the
continuous relaxation share no nonzero recession direction, and they
give an explicit bound under strong convexity and smoothness.
\citet[Theorem~14]{LefebvreSchmidt2025} prove that a finite penalty is
exact at each fixed multiplier under compactness, a differentiable
convex objective, convex slices and a Slater condition on every
equality-feasible slice; the instance of Theorem~\ref{thm:lower}
satisfies these assumptions. Their Theorem~15 computes, in polynomial
time, an infinity-norm penalty that is exact at a given rational
multiplier for MILPs with integer data, and their Table~1 remarks that
the extension to convex MIQPs is easy. For native constraints linear in
all variables and a jointly convex quadratic objective, polynomial
encoding length of an exact penalty is therefore prior work
\citep[Theorem~11]{GuAhmedDey2020}, and so is polynomial-time output of
an exact coefficient for MILPs with integer data
\citep[Theorem~15]{LefebvreSchmidt2025}. Appendix~\ref{app:attainment}
discusses their Example~13: under their supremum-based Definition~1 it
is value-exact at every penalty, but no multiplier attains the dual
value.

\paragraph{Squaring chains and bit complexity.}
The squaring chain and its small residuals are classical;
\citet[Example~6.1]{Bienstock2022} attribute it to Ramana and Khachiyan.
Example~6.2 of the same paper, a nonconvex QCQP with a point of doubly
exponentially small violation, already implies large zero-multiplier
norm penalties. \citet[Sections~2 and~4, Result~4]{Beck2023} use a
bounded convex squaring chain with Slater regularity to show that bilevel
problems are sensitive to feasibility errors. To our knowledge, what is
new in Theorem~\ref{thm:lower} is the exact threshold of the optimized dual for
a compact convex-quadratic model with one binary variable and uniform
strict points, together with its fixed-accuracy consequence. Residuals
of opposite signs on the cost-free slice make the bound hold for every
multiplier.

\paragraph{Real-algebraic tools.}
The upper bounds rest on established tools: radii of balls meeting the
connected components of semialgebraic sets \citep{BasuRoy2010} and
quantitative quantifier elimination \citep{BasuPollackRoy2006,Basu2014}.
Algebraic bounds that depend on the number of quadratic constraints
appear in \citet{GrigorievPasechnik2005} and \citet{KammingaRudolph2025};
Section~\ref{sec:fixedcount} compares them with our argument. To our
knowledge, what is new here are the resulting encoding bounds for exact
penalties, including the dependence on the number of native inequalities
that are nonlinear in the continuous variables while the number of
affine rows is arbitrary.

\paragraph{Hardness of choosing penalties.}
Hardness of choosing penalties, even when the constrained optimum is
known, predates this paper. \citet[Section~IV.A, Lemma~1]{Alessandroni2025}
prove hardness of choosing penalties for exact QUBO reformulations. Their
reduction imposes $x=0$ with a squared residual and a quadratic
objective; on binary variables their penalty $M\sum_ix_i$ is reproduced
by a linear multiplier with augmentation coefficient zero, so it does not
give a positive optimized threshold. \citet[Section~2.2]{Mirkarimi2024}
note that a universally successful, polynomially tuned linear Ising
penalty for knapsack would imply $\mathrm P=\mathrm{NP}$, allowing that no
linear penalty may work on some instances. On the family of
Theorem~\ref{thm:polyfactor}, $D_0=-1<v$, so
no purely linear (Lagrangian) penalty is ever exact and the norm term is
essential. \citet[Definition~1 and Theorem~2]{Alessandroni2026} study
probabilistic feasibility and objective guarantees when exact Gibbs
sampling is used for approximate optimization. Safe conservative
penalties are also known: \citet[Sections~5--6]{GusmeroliWiegele2022}
obtain exact squared penalties for binary quadratic problems with
integer-data linear equality constraints from objective bounds and
residual separation ($\norm{Ax-b}^2\ge1$ by integrality). To our
knowledge, what is new in Theorem~\ref{thm:polyfactor} is the polynomial-factor conclusion for a
norm-penalty threshold that is always positive and is taken after
unrestricted optimization of the multiplier, under the restrictions
above: a binary box, one equality, a linear objective, a known unique
optimizer and a known sufficient coefficient.

\paragraph{Conditioning and perturbation.}
Distance to the boundaries of convex sets is a classical conditioning
device \citep{DunaganSpielmanTeng2011};
\citet[Theorem~3.3 and Corollary~3.4]{BuergisserAmelunxen2012} bound
the relative volume of boundary neighbourhoods of spherically convex
sets within spherical caps.
Theorem~\ref{thm:gridsmoothing} is an explicit application to exact
penalties, with a rational grid, an encoding bound and the infeasibility
alternative; its Gaussian variant (Proposition~\ref{prop:gaussian}) has
the infeasibility alternative but no grid or encoding statement. Their ingredients are known: the boundary-avoidance
principle, the convex repair argument and, for the Gaussian variant, the
tube estimate \citep[Lemma~2.3.4]{DunaganSpielmanTeng2011}. The relation between convex exact
penalties and multiplier bounds is classical
\citep[Theorem~4.2]{FriedlanderTseng2007}, and Lipschitz value bounds
across discrete choices appear in mixed-integer optimal control
\citep[Theorem~3]{GugatHante2017}.

\paragraph{Organization.}
Section~\ref{sec:geometry} defines the exactness notions and the two
thresholds and proves the balanced-mixture, one-equality and slice
results. Section~\ref{sec:lower} proves the lower bound and discusses
which features of the formulation it depends on.
Sections~\ref{sec:upper} and~\ref{sec:fixedcount} prove the encoding
upper bounds in the continuous dimension and in the number of nonlinear
native inequalities. Section~\ref{sec:calibration} proves the
calibration hardness, and Section~\ref{sec:perturbation} the perturbation
results. Section~\ref{sec:discussion} discusses consequences, scope and
open questions. Appendix~\ref{app:attainment} treats Example~13 of
\citet{LefebvreSchmidt2025}, and Appendix~\ref{app:symmetric} gives a
symmetric two-binary variant of the lower bound in which the optimized
and zero-multiplier thresholds coincide.

\section{Exactness and residual geometry}\label{sec:geometry}
Let $X$ be a nonempty compact subset of a finite-dimensional Euclidean
space. Let $f:X\to\R$ and $r:X\to\R^m$ be continuous, and suppose
$F=\{x\in X:r(x)=0\}$ is nonempty. Fix a norm on $\R^m$ and write
\begin{equation}\label{eq:model}
 v=\min_{x\in F}f(x),\qquad
 L_\rho(\lambda)=\min_{x\in X}\{f(x)+\lambda^Tr(x)+\rho\norm{r(x)}\},
 \qquad D_\rho=\sup_{\lambda\in\R^m}L_\rho(\lambda),
\end{equation}
where $\rho\ge0$. The set $X$ retains the \emph{native constraints};
only $r(x)=0$ is relaxed. Weak duality gives $L_\rho(\lambda)\le D_\rho\le v$.
We call $\rho$ \emph{dual-value exact} if $D_\rho=v$, and exact at a
specified multiplier $\lambda$ if $L_\rho(\lambda)=v$. The latter includes
attainment of the dual supremum. We call $(\rho,\lambda)$
\emph{minimizer-set exact} if the minimizers in $L_\rho(\lambda)$ are
exactly $\argmin_F f$. The distinctions are necessary even on compact sets.

Set $\rho_*=\inf\{\rho\ge0:D_\rho=v\}$, with $\inf\varnothing=+\infty$.
For a fixed multiplier, the corresponding value threshold is
\begin{equation}\label{eq:fixed}
 \rho_*(\lambda)=\max\left\{0,\sup_{x\in X\setminus F}
 \frac{v-f(x)-\lambda^Tr(x)}{\norm{r(x)}}\right\}.
\end{equation}
An empty supremum contributes no positive lower bound. This formula follows
by requiring every penalized objective value to be at least $v$.
A finite threshold is exact at that multiplier. Every strictly larger
coefficient is minimizer-set exact: its additional penalty is positive at
each infeasible point. In contrast, finite $\rho_*$ need not give an
attaining multiplier.

Primal convexification descriptions of augmented Lagrangian duals are
established \citep{BolandEberhard2015,FeizollahiAhmedSun2017}.
In particular, \citet[Theorem~1 of the inspected manuscript]{FeizollahiAhmedSun2017}
use a lifted convex hull under a relative-interior feasibility condition.
The following residual-space analogue for compact $X$ needs no
relative-interior condition. Its proof does not assume multiplier
attainment; we make no novelty claim for the convexification argument.

\begin{proposition}[Balanced mixtures]\label{prop:mixture}
For every $\rho\ge0$,
\begin{equation}\label{eq:mixture}
 D_\rho=\min\left\{\sum_{i=1}^p\theta_i
 [f(x_i)+\rho\norm{r(x_i)}]:
 \begin{array}{l}
 1\le p\le m+1,\ x_i\in X,\ \theta_i\ge0,\\
 \sum_i\theta_i=1,\quad\sum_i\theta_i r(x_i)=0
 \end{array}\right\}.
\end{equation}
\end{proposition}
\begin{proof}
The convex hull $K$ of
$\{(r(x),f(x)+\rho\norm{r(x)}):x\in X\}$ is compact.
Let $w=\min\{t:(0,t)\in K\}$. Averaging over a balanced mixture cancels
$\lambda$, so $D_\rho\le w$. For $a<w$, strict separation of $(0,a)$
from $K$ gives $u^Tz+\beta t>\beta a$ for every $(z,t)\in K$.
At $(0,w)$ this implies $\beta>0$. Compactness then gives
$L_\rho(u/\beta)>a$. Letting $a\uparrow w$ proves $D_\rho=w$, without
assuming multiplier attainment.

Carath\'eodory's theorem gives an optimal mixture with at most $m+2$
points. If it has more than $m+1$ positive weights, the columns
$(1,r(x_i))$ are linearly dependent. Both signs of a sufficiently small
weight perturbation along a dependence preserve feasibility.
Optimality forces its objective derivative to be zero. Increase the
perturbation until a weight vanishes; the value is unchanged. This
reduces the support to at most $m+1$.
\end{proof}
The average in \eqref{eq:mixture} contains the norms of the individual
residuals, not the norm of their average, which is zero.
The support bound is sharp: take one point with $(r,f)=(0,0)$ and
$m+1$ points with cost $-1$ and residuals
$e_1,\ldots,e_m,-\sum_j e_j$. At $\rho=0$, an optimal balanced mixture
must give weight $1/(m+1)$ to each of the latter points.

\begin{corollary}\label{cor:ratio}
Let $\mathcal B$ denote the balanced mixtures in \eqref{eq:mixture}, and
write $c_\mu=\sum_i\theta_i f(x_i)$ and
$s_\mu=\sum_i\theta_i\norm{r(x_i)}$. Then
\begin{equation}\label{eq:ratio}
 \rho_*=\max\left\{0,\sup_{\mu\in\mathcal B:s_\mu>0}
                  \frac{v-c_\mu}{s_\mu}\right\}.
\end{equation}
If this number is finite, it is dual-value exact.
\end{corollary}
\begin{proof}
A mixture with $s_\mu=0$ uses only feasible points of positive weight,
so $c_\mu\ge v$. The remaining inequalities in \eqref{eq:mixture} are
$\rho s_\mu\ge v-c_\mu$. They all hold at their finite supremum.
\end{proof}

\begin{proposition}[One relaxed equality]\label{prop:scalar}
Suppose $m=1$, use absolute value, and assume residuals of both signs occur.
Define
\[
 A_+=\sup_{r(x)>0}\frac{v-f(x)}{r(x)},\qquad
 A_-=\sup_{r(x)<0}\frac{v-f(x)}{-r(x)}.
\]
If both are finite, then
\begin{equation}\label{eq:scalar}
 \rho_*=\max\{0,(A_++A_-)/2\},\qquad
 L_\rho(\lambda)=v\ \Longleftrightarrow\
 A_+-\rho\le\lambda\le\rho-A_-.
\end{equation}
If either is infinite, no finite penalty is dual-value exact.
At multiplier zero the threshold is $\max\{0,A_+,A_-\}$.
\end{proposition}
\begin{proof}
The pointwise inequalities for positive and negative residuals give the
interval in \eqref{eq:scalar}. To see that an exact supremum cannot avoid
this interval, fix one point of each residual sign. Evaluating any sequence
$L_\rho(\lambda_j)\to v$ at these two points bounds $\lambda_j$ above
and below. Along a convergent subsequence, upper semicontinuity of
$L_\rho$ gives an attaining multiplier. Hence exactness requires and is
implied by a nonempty interval, giving all the claims.
\end{proof}
If residuals occur on only one side of zero, every balanced mixture is
supported on $F$, so $D_\rho=v$ for all $\rho\ge0$. Attainment can fail:
for $X=[0,1]$, $r(x)=x$, and $f(x)=-\sqrt{x}$, every finite multiplier
and penalty give $L_\rho(\lambda)<0=v$. Appendix~\ref{app:attainment}
records the same distinction in a mixed-integer quadratic example from
the literature.

\begin{lemma}[A bound from integer slices]\label{lem:slices}
Suppose $X$ is the union of finitely many nonempty compact slices $X_z$,
$f\ge L$ on $X$, and $v\le U$. On every slice meeting $F$, suppose
there is $\lambda_z$ such that
\[
 f(x)+\lambda_z^Tr(x)\ge v_z:=\min_{X_z\cap F}f,
 \qquad \norm{\lambda_z}_*\le M\quad(x\in X_z),
\]
where $\norm{\cdot}_*$ is the dual norm. Suppose every other slice has
$\min_{X_z}\norm{r}\ge\delta>0$. Then
\begin{equation}\label{eq:slicebound}
 \rho\ge\max\{M,(U-L)/\delta\}
\end{equation}
is exact at multiplier zero. Any strictly larger coefficient is
minimizer-set exact. If all slices meet $F$, omit the second term.
\end{lemma}
\begin{proof}
On a feasible slice, the dual-norm inequality gives
$f+\rho\norm r\ge f+\lambda_z^Tr\ge v_z\ge v$.
On an infeasible slice, $f+\rho\norm r\ge L+\rho\delta\ge U\ge v$.
A primal minimizer attains $v$. Strictly increasing $\rho$ excludes every
infeasible point from the minimizing set.
\end{proof}
Thus feasible-slice multiplier bounds and infeasible-slice separation
control different parts of the penalty. For convex slices, Slater regularity can supply
the former without giving a useful quantitative bound on the latter.

\section{An exponential encoding lower bound}\label{sec:lower}
All encoding statements use explicit rational coefficients and ordinary
binary representations of their numerators and positive denominators.
A sparse encoding records variable indices as well as nonzero coefficients.
The norm and the relaxed equality are fixed as displayed.

For $n\ge1$, let
\begin{equation}\label{eq:chain}
 A_n=\{a\in[0,1]^n:a_1\ge1/4,\ a_{i+1}\ge a_i^2\ (1\le i<n)\},
 \qquad \delta_n=2^{-2^n}.
\end{equation}
Induction shows $a_i\ge2^{-2^i}$, with all bounds attained by
$\bar a_i=2^{-2^i}$. Define
\begin{equation}\label{eq:onebinary}
 X_n=\{(a,y,q):a\in A_n,\ -1\le y\le1,\ q\in\{0,1\},\
                  y\ge a_n-2(1-q)\}.
\end{equation}
We minimize $-q$ subject to $(a,y,q)\in X_n$ and $y=0$.
In \eqref{eq:model}, take $f=-q$, $r=y$, and $\norm r=|y|$.

\begin{theorem}\label{thm:lower}
The problem defined using $X_n$ in \eqref{eq:onebinary} has value $v_n=0$, and for every $\rho\ge0$,
\begin{equation}\label{eq:dualformula}
 D_{n,\rho}=\min\left\{0,\frac{2\rho\delta_n-1}{1+\delta_n}\right\}.
\end{equation}
Its dual supremum is attained and its least dual-value exact penalty is
\begin{equation}\label{eq:lowerthreshold}
 \rho_{*,n}=\frac1{2\delta_n}=2^{2^n-1}.
\end{equation}
The instance has one binary variable, $n+1$ continuous variables, a linear
objective, and convex quadratic native constraints with coefficients from
a fixed finite set. Its sparse encoding length is $O(n\log(n+1))$, whereas
every rational exact penalty needs at least $2^n$ numerator bits.
Every equality-feasible integer slice and every native integer slice have
strict points whose continuous inequality slacks are at least $1/8$.
\end{theorem}
\begin{proof}
Projection onto $(q,y)$ gives exactly
\begin{equation}\label{eq:projection}
 (\{0\}\times[-1,1])\ \cup\ (\{1\}\times[\delta_n,1]).
\end{equation}
The inclusion follows from the chain; the reverse follows by setting
$a=\bar a$. At $q=0$ the lower bound $\delta_n-2$ is below $-1$.
Consequently $y=0$ forces $q=0$ and $v_n=0$.

For any multiplier, the points $(q,y)=(0,-1)$ and $(1,\delta_n)$
give values $\rho-\lambda$ and $-1+(\rho+\lambda)\delta_n$.
Their weighted average, with weights $\delta_n/(1+\delta_n)$ and
$1/(1+\delta_n)$, is $(2\rho\delta_n-1)/(1+\delta_n)$.
The point $(0,0)$ gives zero. These witnesses prove the upper bound
in \eqref{eq:dualformula} for every multiplier.

For $0\le\rho\le1/(2\delta_n)$, take
\[
 \lambda_\rho=\frac{1+\rho(1-\delta_n)}{1+\delta_n}\ge\rho.
\]
On the first interval of \eqref{eq:projection}, the minimum occurs at
$y=-1$; on the second it occurs at $y=\delta_n$. Their values coincide
with the proposed dual value. For $\rho\ge1/(2\delta_n)$, take
$\lambda=\rho$. The penalized objective is nonnegative on both intervals
and zero at $(0,0)$. This proves the formula and attainment.

Each inequality $a_i^2-a_{i+1}\le0$ is convex; all other constraints
are affine. Across all constraints there are $O(n)$ nonzero coefficients,
each with variable indices of $O(\log(n+1))$ bits.
The integer $2^{2^n-1}$ has exactly $2^n$ bits.
If a rational $p/d$ with positive integers $p,d$ is at least this value,
then $p\ge2^{2^n-1}$ as well.

For the equality-feasible slice $q=0$, set $a_i=3/8$ and $y=0$.
For the native slice $q=1$, set $a_i=3/8$ and $y=1/2$.
The seed inequality has slack $1/8$, each chain inequality has slack
$15/64$, and the variable bounds and linking inequality have slack at
least $1/8$. The native $q=0$ slice uses the first point. These are
uniform strict points, although the equality-infeasible $q=1$ slice
comes within $\delta_n$ of $y=0$.
\end{proof}
The construction satisfies Assumptions~1--5 of
\citet{LefebvreSchmidt2025}, so their Theorem~14 ensures a finite exact
penalty at every fixed multiplier. Theorem~\ref{thm:lower} concerns the
encoding length of such penalties, even after optimizing the multiplier.

At $\rho=\rho_{*,n}$, Proposition~\ref{prop:scalar} gives the unique exact
multiplier $\lambda=\rho_{*,n}$: here $A_+=1/\delta_n$ and $A_-=0$.
Both witnesses in the proof then tie with feasible optima.
For every $\rho>\rho_{*,n}$, the fixed multiplier
$\lambda=\rho_{*,n}$ makes all infeasible objective values positive,
so the minimizer set is exact. The zero-multiplier value threshold is
$1/\delta_n$, twice the optimized threshold.

\begin{corollary}[Fixed additive accuracy]\label{cor:accuracy}
For $0\le\epsilon<1$, the condition $v_n-D_{n,\rho}\le\epsilon$ is
equivalent to
\begin{equation}\label{eq:accuracy}
 \rho\ge\max\left\{0,\frac{1-\epsilon(1+\delta_n)}{2\delta_n}\right\}.
\end{equation}
For each fixed $\epsilon<1$ and all sufficiently large $n$, such a rational
penalty needs at least $2^n-O_\epsilon(1)$ numerator bits.
\end{corollary}
\begin{proof}
The first claim follows from \eqref{eq:dualformula}.
For large $n$, $\epsilon\delta_n\le(1-\epsilon)/2$, so the right-hand side
is at least $(1-\epsilon)/(4\delta_n)$. Taking binary logarithms gives
the bit bound.
\end{proof}
In particular, accuracy $\epsilon=1/2$ requires
$\rho\ge(1-\delta_n)/(4\delta_n)$. The objective range is one, so this
is a fixed accuracy relative to that normalization. The lower bound is
exponential in continuous dimension and superpolynomial in input length;
it is not a claim of an exponential bound in the full sparse input length.
Nor is it a running-time lower bound for solving the primal, whose value
is immediate from \eqref{eq:projection}.

\subsection{A symmetric variant and the role of the penalty function}
A two-binary companion makes multiplier cancellation symmetric.
Replace the final inequality in \eqref{eq:onebinary} by
\begin{equation}\label{eq:symmetric}
 y\ge a_n-(1-q)-2(1-b),\qquad
 y\le-a_n+(1-q)+2b,\qquad b\in\{0,1\}.
\end{equation}
Keep the same objective and relaxed equality. Projection now gives
\[
 (\{0\}\times[-1,1])\ \cup\
 (\{1\}\times([-1,-\delta_n]\cup[\delta_n,1])).
\]
Indeed, at $a=\bar a$ the two $q=0$ intervals are
$[-1,1-\delta_n]$ and $[\delta_n-1,1]$, which cover $[-1,1]$;
at $q=1$ they are the two outer intervals. If $c=\rho-|\lambda|$,
minimizing on these intervals yields
\begin{equation}\label{eq:symdual}
 L_{n,\rho}^{\rm sym}(\lambda)=\min\{0,-1+c\delta_n,-1+c\},\qquad
 D_{n,\rho}^{\rm sym}=\min\{0,-1+\rho\delta_n\}.
\end{equation}
For the first formula, the $q=0$ branch has value $\min(0,c)$ and
is dominated by one of the three displayed terms. The second follows
because $c\le\rho$ and $\lambda=0$ attains the bound.
Thus the least exact penalty is $\delta_n^{-1}$, and gap at most
$\epsilon<1$ requires exactly $\rho\ge(1-\epsilon)\delta_n^{-1}$.
Any strictly larger penalty with $\lambda=0$ is minimizer-set exact.
Strict points again use $a_i=3/8$, with $y=0$ for $q=0$ and
$y=(2b-1)/2$ for $q=1$; their minimum continuous inequality slack is $1/8$.

The fixed norm is essential. Let $p_n=2^{-n}$, a rational exponent with
$O(n)$ bits. Then $\delta_n^{p_n}=1/2$. For either construction,
$f+2|y|^{p_n}$ has minimum zero: at $q=0$ it is nonnegative and at
$q=1$ it is at least $-1+2\delta_n^{p_n}=0$.
Any coefficient larger than two gives minimizer-set exactness.
This instance-dependent augmentation is not Lipschitz at zero and is
not a norm. It changes the encoding question and the continuous
minimization problem. Likewise, succinct exponent expressions or a
rescaled linking equality change the representation or formulation to
which Theorem~\ref{thm:lower} applies.

\section{An encoding upper bound in continuous dimension}\label{sec:upper}
We now consider rational quadratic input of total binary length $N\ge2$:
\begin{equation}\label{eq:quadraticmodel}
 v=\min\{f(x,z):(x,z)\in X,\ Ax+Bz-b=0\},\qquad
 x\in\R^n,\ z\in\Z^p,
\end{equation}
where $n\ge1$ and the residual $r(x,z)=Ax+Bz-b$ has $m\ge1$
coordinates. The native set is
specified by finite rational coordinate bounds and inequalities
$g_i(x,z)\le0$, all of total degree at most two; $f$ has the same degree
bound. Let $s$ count all native inequalities, including the boxes. Affine
equalities may be represented by paired inequalities. The input lists
expanded polynomials, with coefficients and variable indices encoded
explicitly; in particular, $s,m,n\le N$.

For every integer assignment, assume that $f(\cdot,z)$ and each
$g_i(\cdot,z)$ are convex on $\R^n$. Write
$X_z=\{x:(x,z)\in X\}$. A slice is \emph{equality-feasible} if
$X_z\cap\{x:Ax+Bz=b\}\ne\varnothing$.
Suppose the original feasible set is nonempty and, on every such slice,
this intersection contains a point strictly satisfying every native
inequality nonlinear in $x$; native affine inequalities may hold weakly.
This is the \emph{refined Slater condition}. No margin is given,
and no Slater condition is imposed on equality-infeasible slices.
Empty native slices may be discarded. The finite boxes leave only finitely
many integer assignments.

Write $\operatorname{bit}(a)=1+\lfloor\log_2 a\rfloor$ for a positive
integer $a$. All constants implicit below are absolute and independent of
the integer dimension and the number of slices.

\begin{theorem}\label{thm:generalupper}
Under these assumptions there is an integer $\rho\ge1$ with
\begin{equation}\label{eq:generalbits}
 \operatorname{bit}(\rho)\le (N+1)2^{O(n)}
\end{equation}
such that $(\rho,0)$ is minimizer-set exact for either
$\norm{\cdot}_\infty$ or $\norm{\cdot}_1$.
\end{theorem}
We use the following consequence of effective bounds for real algebraic
sets. A \emph{basic closed set} is a conjunction of polynomial equalities
and weak inequalities.

\begin{lemma}[Effective radii]\label{lem:radii}
Let at most $S$ integer polynomials in $d$ variables have degree at most
two and coefficient bitsize at most $\tau\ge1$. There is a radius $R\ge1$
with
\begin{equation}\label{eq:radius}
 \log_2 R\le(\tau+\log(S+1)+1)2^{O(d)}
\end{equation}
that meets every nonempty basic closed set defined by these polynomials
and contains every such set that is bounded. The same statements hold
for a finite union of such sets.
\end{lemma}
\begin{proof}
Apply the component-meeting and bounded-component radii of
\citet[Theorems~3--4]{BasuRoy2010} to weak sign conditions and take the
larger radius. Their explicit formulas give \eqref{eq:radius} for fixed
degree. A conjunction that leaves some polynomial signs unspecified is
a finite union of compatible weak sign conditions. For a union, select
one nonempty member to obtain a point; if the union is bounded, every
member and every component is bounded. The source proof's factorial-bit
term, retained when majorizing its meeting radius, is also absorbed by
$2^{O(d)}$.
\end{proof}
Only closed conjunctions and their finite unions are used here; the lemma
makes no claim about arbitrary strict sign conditions.

\subsection{Uniform coefficient bounds}
The explicit input gives coordinate bounds of magnitude $2^{O(N)}$ and
$O(N)$ encoded monomials. Multiplying by the product of all input
denominators clears coefficients with $O(N)$ bits. Substituting any boxed
integer assignment into a degree-two monomial adds only $O(N)$ bits;
summing monomials adds $O(\log N)$ more. Differentiation and introducing
multiplier variables preserve this uniform $O(N)$ coefficient-height
bound. Thus every degree-two polynomial system below can be cleared with
\begin{equation}\label{eq:height}
 \tau=O(N).
\end{equation}
We clear the polynomial equations of certificate systems, not the
residual in the penalty; the multiplier coordinates keep their original
meaning. A box bound $|f|\le F_0$, with an integer $F_0\ge1$ and
$\operatorname{bit}(F_0)=O(N)$, is equally immediate. For example, if
$H\ge1$ bounds all coordinate magnitudes and $C$ is the sum of the
absolute objective coefficients, take $F_0=1+\lceil CH^2\rceil$.

\subsection{Feasible slices: a sparse multiplier certificate}
Fix an equality-feasible assignment and suppress $z$. Refined Slater
ensures a convex KKT certificate $(x^*,\mu,\lambda)$, with
$\mu\ge0$. To apply the usual relative-interior form of Slater, one may
blend the strictly feasible point with a point in the relative interior
of the polyhedron defined by the affine constraints. For a sufficiently
small blend, all nonlinear inequalities remain strict. The resulting
convex optimality conditions express the normal of that polyhedron by
its active affine rows, giving the stated KKT certificate even when
these rows are dependent.

Let $e=\operatorname{rank} A$ and select independent rows $J$ spanning
its row space. Project stationarity onto the quotient by this row space.
The projected vector $-\nabla f(x^*)$ is in the cone generated by active
native gradients, in dimension $n-e$. Conic Carath\'eodory therefore
selects at most $n-e$ gradients, indexed by $I$. The remaining vector
is in the row space of $A$. Hence $|I|+|J|\le n$, and the following
system has a solution:
\begin{equation}\label{eq:sparsekkt}
 \begin{aligned}
 &g_i(x)\le0\quad(1\le i\le s),\qquad Ax+Bz-b=0,\\
 &\mu_i\ge0,\quad g_i(x)=0\quad(i\in I),\\
 &\nabla f(x)+\sum_{i\in I}\mu_i\nabla g_i(x)
                  +\sum_{j\in J}\lambda_j A_j^T=0.
 \end{aligned}
\end{equation}
This is a basic closed set in at most $2n$ variables, defined by at most
$s+m+2n$ degree-two polynomials: the active equations reuse the native
polynomials $g_i$. Replacing complementarity products by these equations
keeps degree two. The selected original
certificate proves nonemptiness; every point in the new system remains
a KKT certificate.

By Lemma~\ref{lem:radii}, some certificate has all coordinates bounded
by $R_M$, where
\begin{equation}\label{eq:RM}
 \log_2 R_M\le (\tau+\log(s+m+2n+1)+1)2^{O(n)}.
\end{equation}
Set omitted linking multipliers to zero. The resulting multiplier has
$\norm{\lambda_z}_1\le nR_M$. Convexity, stationarity and
complementarity imply, throughout the native slice,
\begin{equation}\label{eq:slicemultiplier}
 f(x,z)+\lambda_z^T(Ax+Bz-b)\ge v_z.
\end{equation}
Indeed, the full convex Lagrangian is minimized at the KKT point and
has value $v_z$ there; deleting its nonpositive native-inequality terms
on $X_z$ only increases it. We bound one suitable certificate, not all
optimal multipliers, whose set can be unbounded.

\subsection{Infeasible slices: a reciprocal graph}
For a nonempty equality-infeasible native slice, compactness gives
$\delta_z=\min_{x\in X_z}\norm{Ax+Bz-b}_\infty>0$.
Write $r_j=r_j(x,z)$ and consider
\begin{equation}\label{eq:reciprocal}
 T_z=\left\{(x,t):\begin{array}{l}
 x\in X_z,\quad t\ge0,\quad -1\le tr_j\le1\ (1\le j\le m),\\
 tr_j=1\ \text{or}\ tr_j=-1\ \text{for at least one }j
 \end{array}\right\}.
\end{equation}
This is exactly the graph of $t=1/\norm{r(x,z)}_\infty$. It is
nonempty and compact, since the denominator is bounded away from zero
on $X_z$. Selecting a coordinate and a sign in the last line expresses
it as a union of $2m$ basic closed sets. The equality uses one of the
polynomials already present in $-1\le tr_j\le1$, so only
$s+2m+1$ degree-two polynomials in $n+1$ variables are needed.
The bounded-set part of Lemma~\ref{lem:radii} gives
\begin{equation}\label{eq:RD}
 \delta_z^{-1}\le R_D,\qquad
 \log_2 R_D\le(\tau+\log(s+2m+2)+1)2^{O(n)}.
\end{equation}
No KKT condition or rationality of $\delta_z$ is used on these slices.

\begin{proof}[Proof of Theorem~\ref{thm:generalupper}]
Choose an integer $E\ge0$ with $2^E\ge\max\{R_M,R_D\}$ satisfying
the preceding bounds, and set
\begin{equation}\label{eq:sufficientrho}
 \rho=(n+2F_0+1)2^E.
\end{equation}
This is strictly larger than both $nR_M$ and $2F_0R_D$.
Lemma~\ref{lem:slices}, with objective bounds $-F_0,F_0$, proves
minimizer-set exactness for the infinity norm. In detail, on a feasible
slice a point of nonzero residual has penalized value at least
$v_z+(\rho-nR_M)\norm r_\infty>v$, while on an infeasible slice it
has value at least $-F_0+\rho/R_D>F_0\ge v$.
The same coefficient works for the one-norm because
$\norm r_1\ge\norm r_\infty$ and both vanish on the original feasible
set. Equations \eqref{eq:height}--\eqref{eq:RD} and
$\operatorname{bit}(F_0)=O(N)$ give \eqref{eq:generalbits}. All bounds are
uniform over boxed integer assignments, so no factor counting those
assignments appears.
\end{proof}
Together with Theorem~\ref{thm:lower}, this gives exponential upper
and lower dependence of penalty bit length on continuous dimension.
The constants in the exponents and the dependence on the full sparse
input length are not asserted to be sharp.

\section{Polynomial encoding with few nonlinear inequalities}
\label{sec:fixedcount}
This section proves a second upper bound: if at most $k$ native
inequalities are nonlinear in the continuous variables, some sufficient
coefficient has $N^{O(k+1)}$ bits. For fixed $k$ this is polynomial in the
input length, however large the continuous dimension, the number of affine
rows or the number of integer assignments. In practical terms, many linear
rows together with a few convex quadratic rows cannot produce the doubly
exponential growth of Section~\ref{sec:lower}.

We work in the setting of Section~\ref{sec:upper}. Let $k$ be an upper
bound on the number of native inequalities $g_i(x,z)\le0$ that are
nonlinear in the continuous variables $x$. Constraints that are affine in
$x$ are not counted, even if they contain integer variables, such as
$x_1z_1+z_1^2\le3$; the objective is not counted either. As noted after
Assumption~\ref{ass:model}, whether a row is nonlinear in $x$ does not
depend on $z$.

\begin{theorem}\label{thm:fixedcount}
Suppose that Assumption~\ref{ass:model} holds and that at most $k$
native inequalities are nonlinear in the continuous variables. Then there is an integer $\rho\ge1$ with
\begin{equation}\label{eq:fixedbits}
 \operatorname{bit}(\rho)\le N^{O(k+1)}
\end{equation}
that is minimizer-exact at multiplier zero for both the infinity norm and
the one-norm; the same $\rho$ works for both. The constant in the exponent
is absolute.
\end{theorem}
Neither Theorem~\ref{thm:generalupper} nor Theorem~\ref{thm:fixedcount}
implies the other: the first is polynomial for fixed $n$, the second for
fixed $k$.

\paragraph{Proof strategy.}
In the proof of Theorem~\ref{thm:generalupper}, both the separation of
equality-infeasible slices and the linking multipliers of
equality-feasible slices are bounded through radii of semialgebraic sets
in about $2n$ variables, which costs a factor $2^{O(n)}$ in the exponent.
Here each slice needs only one scalar. For an equality-infeasible slice
it is the separation $\delta_z=\min_{X_z}\norm r_\infty$. For an
equality-feasible slice it is a Slater margin $\gamma_z$, the largest
amount (capped at one) by which some point of the slice with $r=0$
satisfies all nonlinear rows strictly. Both are optimal values of convex quadratic problems with
at most $k$ nonlinear rows. Lemma~\ref{lem:nonzerovalue} below bounds such
nonzero optimal values below by $2^{-N^{O(k+1)}}$; here $k$ enters only
the exponent of $N$, and $n$ does not enter it.
The Slater margin then bounds the multipliers of the nonlinear rows, and
the linking multipliers are recovered from the remaining affine
stationarity equations by linear algebra (Cramer's rule).

\subsection{An algebraic bound on optimal values}
The lemma says that the optimal value of a convex quadratic problem with
$k$ quadratic rows is an algebraic number whose degree and coefficient
size are polynomial in the input for fixed $k$; a nonzero such number
cannot be smaller than $2^{-\mathrm{poly}}$.

\begin{lemma}\label{lem:nonzerovalue}
Minimize a rational convex quadratic function over a nonempty set
described by a rational polyhedron with finite coordinate bounds and at
most $k$ rational convex quadratic inequalities. Let $N\ge2$ be the input
length of this problem, encoded explicitly as in
Assumption~\ref{ass:model}, and let $\omega$ be its optimal value. Then
$\omega$ is a root of a nonzero univariate integer polynomial whose degree
and coefficient bitsize are at most $N^{O(k+1)}$. In particular, if
$\omega\ne0$, then
\begin{equation}\label{eq:nonzero}
 |\omega|\ge2^{-N^{O(k+1)}}.
\end{equation}
No Slater condition or other constraint qualification is required.
\end{lemma}
\begin{proof}
The proof has three steps. (1) We delete the affine inequalities that are
inactive at an optimum and substitute the active affine face; this leaves
an unconstrained affine parametrization and at most $k+1$ quadratic rows.
(2) We regularize the objective and relax the rows by a parameter
$\eta>0$, so that KKT conditions hold with a positive definite matrix, and
eliminate the primal variables with the adjugate; this gives polynomial
conditions in $\eta$, the value $w$ and the multipliers. (3) We express
the limit $\eta\downarrow0$ by a formula with two quantifier blocks, of
sizes $1$ and at most $k+3$, and apply effective quantifier elimination,
whose degree and bitsize bounds are exponential only in the block sizes.

Write the problem in variables $x\in\R^n$ as the minimization of
$\phi_0(x)$ over the points of the polyhedron $P$ that satisfy
$\phi_i(x)\le0$ for $1\le i\le k$. The feasible set is nonempty and
compact, so an optimum $x^*$ exists. Note that $n\le N$.

\emph{Step 1: affine-face reduction.}
Let $E$ be the affine space defined by all equalities of $P$ and all
affine inequalities of $P$ that are active at $x^*$. The remaining affine
inequalities of $P$ have positive slack at $x^*$, hence on a neighborhood
$W$ of $x^*$. Deleting them leaves $x^*$ optimal over
$E\cap\{\phi_i\le0,\ 1\le i\le k\}$. Indeed, suppose a point $y$ of this
set had $\phi_0(y)<\phi_0(x^*)$. By convexity, every point
$x^*+\theta(y-x^*)$ with $0<\theta\le1$ lies in $E$, satisfies the
inequalities $\phi_i\le0$, and has objective value below $\phi_0(x^*)$.
For small $\theta$
it also lies in $W$, so it satisfies the deleted rows and is feasible for
the original problem. This contradicts the optimality of $x^*$.

Choose independent equations defining $E$. Gaussian elimination writes
$E=\{x_0+Vu:u\in\R^{n'}\}$, where $x_0$ and $V$ are rational with bit
length polynomial in $N$, and the free coordinates $u$ are $n'$ of the
coordinates of $x$; thus $n'\le n\le N$. If $n'=0$, then $E=\{x^*\}$ with
$x^*=x_0$, so $\omega=\phi_0(x_0)$ is rational with bit length polynomial
in $N$; it is the root of a linear integer polynomial with such
coefficients, and the lemma follows. Assume $n'\ge1$. Let
$R_{\mathrm{box}}\ge1$ be an integer of bit length polynomial in $N$ that
bounds the absolute values of all coordinates on $P$. Since $u^*$ is a
subvector of $x^*$, we have $\norm{u^*}_2^2\le n'R_{\mathrm{box}}^2$. Add
the ball constraint $\norm u_2^2-n'R_{\mathrm{box}}^2-1\le0$. It holds
strictly at $u^*$, keeps the optimal value $\omega$ attained at $u^*$, and
makes the feasible set compact.

We now have an objective $\psi_0(u)=\phi_0(x_0+Vu)$ and $k'\le k+1$
convex quadratic inequalities $\psi_i(u)\le0$: the substituted rows
$\phi_i(x_0+Vu)\le0$ and the ball. There are no remaining affine
constraints, and all coefficients have bitsize polynomial in $N$. Some
$\psi_i$ may be affine or constant; this does not affect the argument.
Write
\[
 \psi_i(u)=\tfrac12u^TQ_iu+a_i^Tu+c_i\qquad(0\le i\le k'),
\]
with symmetric $Q_i\succeq0$. The matrices $Q_i$ may be singular.

\emph{Step 2: regularization and elimination.}
For $0<\eta<1$, define
\begin{equation}\label{eq:regularized}
 \omega(\eta)=\min\{\psi_0(u)+\eta\norm u_2^2:
                         \psi_i(u)\le\eta\ (1\le i\le k')\}.
\end{equation}
The relaxed ball bounds every feasible point, so the feasible set is
compact, and it contains $u^*$; hence an optimum exists. The point $u^*$
satisfies every relaxed inequality strictly, so it is a Slater point of
this convex problem, and nonnegative KKT multipliers $\mu_i$ exist at
every optimum. (Strict convexity of the objective only adds uniqueness of
the optimum, which we do not need.) Stationarity reads
\begin{equation}\label{eq:matrixstationarity}
 H(\eta,\mu)\,u=-a_0-\sum_{i=1}^{k'}\mu_i a_i,\qquad
 H(\eta,\mu)=Q_0+2\eta I+\sum_{i=1}^{k'}\mu_iQ_i.
\end{equation}
The term $2\eta I$ makes $H(\eta,\mu)$ positive definite for all $\eta>0$
and all $\mu\ge0$, even when every $Q_i$ is singular. This allows us to
solve for $u$. Let
\[
 \Delta(\eta,\mu)=\det H(\eta,\mu),\qquad
 U(\eta,\mu)=-\operatorname{adj}\bigl(H(\eta,\mu)\bigr)
             \Bigl(a_0+\sum_{i=1}^{k'}\mu_ia_i\Bigr).
\]
The entries of $H(\eta,\mu)$ are affine in $(\eta,\mu)$, so $\Delta$ and
the entries of $U$ are polynomials of degree at most $n'$, and
$u=U/\Delta$. We abbreviate $\Delta=\Delta(\eta,\mu)$ and $U=U(\eta,\mu)$
and define
\begin{align}
 \Phi_i&=\tfrac12U^TQ_iU+\Delta\,a_i^TU
                  +(c_i-\eta)\Delta^2\quad(1\le i\le k'),
                  \label{eq:Gconstraint}\\
 \Phi_0&=\tfrac12U^TQ_0U+\Delta\,a_0^TU
                  +(c_0-w)\Delta^2+\eta\,U^TU.
                  \label{eq:Gobjective}
\end{align}
For $\Delta\ne0$ these are
$\Phi_i=\Delta^2(\psi_i(U/\Delta)-\eta)$ and
$\Phi_0=\Delta^2(\psi_0(U/\Delta)+\eta\norm{U/\Delta}_2^2-w)$.
Consider the polynomial system
\begin{equation}\label{eq:Ksystem}
 \Sigma(\eta,w,\mu):\quad
 \Delta>0,\quad
 \mu_i\ge0,\ \Phi_i\le0,\ \mu_i\Phi_i=0\ (1\le i\le k'),
 \quad \Phi_0=0.
\end{equation}
The condition $\Delta>0$ is implied by $\eta>0$ and $\mu\ge0$; we keep it
for clarity. For $0<\eta<1$, the system $\Sigma(\eta,w,\mu)$ has a
solution $\mu$ exactly when $w=\omega(\eta)$. If $w=\omega(\eta)$, an
optimum and its KKT multipliers give a solution: stationarity gives
$u=U/\Delta$, and feasibility, complementarity and the optimal value give
the remaining conditions. Conversely, if $\mu$ solves the system, then
$u=U/\Delta$ satisfies stationarity, feasibility and complementarity, so
$(u,\mu)$ is a KKT pair of a convex problem; hence $u$ is optimal and
$\Phi_0=0$ gives $w=\omega(\eta)$.

Every polynomial in \eqref{eq:Ksystem} has degree at most $2n'+2$. Its
coefficients have bitsize polynomial in $N$. To see this, clear the
denominators of the entries of the matrices $Q_i$ by one positive common
denominator. Each coefficient of the determinant, as a polynomial in
$(\eta,\mu)$, is then a sum of at most $n'!\,(k'+2)^{n'}$ products of at
most $n'$ integers of polynomial bitsize; since $n'\le N$ and $k'\le N+1$,
its bitsize is polynomial in $N$. The adjugate and the polynomials
$\Phi_i$ have the same property. This bounds coefficient heights even when
the number of monomials is large. Multiplying each polynomial by a
positive integer that clears its remaining denominators does not change
the sign conditions.

\emph{Step 3: the limit and its algebraic height.}
We have $\omega(\eta)\to\omega$ as $\eta\downarrow0$. For the upper
limit, $u^*$ is feasible for \eqref{eq:regularized}, so
$\omega(\eta)\le\omega+\eta\norm{u^*}_2^2$. For the lower limit, take any
sequence $\eta_j\downarrow0$ and regularized optima $u_j$. They lie in a
common ball, so a subsequence converges to some $\bar u$. The limit
satisfies $\psi_i(\bar u)\le0$ for all $i$, so
$\psi_0(\bar u)\ge\omega$, and $\omega(\eta_j)\ge\psi_0(u_j)$ gives
$\liminf_j\omega(\eta_j)\ge\omega$ along the subsequence. Applying this to
a subsequence that realizes the $\liminf$ proves the claim. No bound on
the regularized multipliers is needed.

Consequently the following formula, with one free variable $\alpha$,
defines exactly the set $\{\omega\}$:
\begin{equation}\label{eq:limitformula}
 \forall\xi\ \exists(\eta,w,\mu)\ \left[\ \xi\le0\ \lor\
   \left(
   \begin{array}{l}
    0<\eta<1,\quad\eta<\xi,\\
    -\xi<w-\alpha<\xi,\quad\Sigma(\eta,w,\mu)
   \end{array}\right)\right].
\end{equation}
Since $\xi\le0$ does not involve $\eta$, $w$ or $\mu$, this prenex
formula is equivalent to
$\forall\xi\,[\xi\le0\lor\exists(\eta,w,\mu)(\ldots)]$. For
$\alpha=\omega$, convergence gives, for every $\xi>0$, some
$0<\eta<\min\{1,\xi\}$ with $|\omega(\eta)-\omega|<\xi$, and
$w=\omega(\eta)$ with KKT multipliers $\mu$ is a witness. Conversely, if
the formula holds for $\alpha$, then witnesses for $\xi=1/j$ give
$\eta_j<1/j$ and $w_j=\omega(\eta_j)$ with $|w_j-\alpha|<1/j$; since
$\omega(\eta_j)\to\omega$, we get $\alpha=\omega$.

We apply effective block quantifier elimination
\citep[Theorem~14.16]{BasuPollackRoy2006}, also stated in
\citet[Theorem~2.27]{Basu2014}. For a prenex formula with quantifier
blocks of sizes $b_1$ and $b_2$, $\ell$ free variables, polynomials of
degree at most $D$, and integer coefficients of bitsize at most $\tau$, it
gives an equivalent quantifier-free formula whose polynomials have degree
at most $D^{O(b_1)O(b_2)}$ and integer coefficients of bitsize at most
$\tau D^{O(b_1)O(b_2)O(\ell)}$. We use this bound on coefficient heights,
not only the bound on the number of arithmetic operations. For
\eqref{eq:limitformula} the parameters are
\[
 b_1=1\ (\text{the variable }\xi),\qquad
 b_2=k'+2\le k+3\ (\text{the variables }\eta,w,\mu),\qquad
 \ell=1\ (\text{the variable }\alpha),
\]
\[
 D=2n'+2=O(n'+1),\qquad \tau=N^{O(1)}.
\]
Since $n'\le N$, we have $D^{O(k+1)}=N^{O(k+1)}$. The quantifier-free
formula therefore involves univariate integer polynomials in $\alpha$ of
degree and coefficient bitsize $N^{O(k+1)}$.

Remove the identically zero polynomials from this formula; their signs
are constant. At least one remaining polynomial vanishes at $\omega$.
Otherwise every remaining polynomial would have a constant nonzero sign on
a neighborhood of $\omega$, so the formula would have a constant truth
value there, whereas it defines the singleton $\{\omega\}$. This proves
the claim on degree and height.

Now let $\omega\ne0$. Divide that polynomial by the largest power of the
indeterminate that divides it. The quotient still vanishes at $\omega$,
has a nonzero constant term, and its coefficients are among those of the
original polynomial; let $\tau_1$ bound their bitsize. The reciprocal
polynomial has the root $1/\omega$ and a leading coefficient of absolute
value at least one, so the Cauchy bound gives
$|1/\omega|\le1+2^{\tau_1}$, that is, $|\omega|\ge1/(1+2^{\tau_1})$.
Absorbing the additive constant in the exponent proves \eqref{eq:nonzero}.
\end{proof}

The argument neither finds the active affine face algorithmically nor
needs a genericity assumption. Algebraic methods for few quadratic
constraints have direct precedents.
\citet[Theorem~1.2]{GrigorievPasechnik2005} give sampling and coefficient
bounds for zero sets of polynomials in quadratic maps; their Theorem~1.5
states an
optimization version, whose proof is deferred to a continuation of their
paper. \citet[Sections~6.2--7.2 and~8.5]{KammingaRudolph2025} give
reduced-variable formulas for optimization under a boundedness hypothesis
on the solution set, and an application to QCQP with a fixed total number
of constraints. After our affine-face reduction, squared slack variables
would turn the inequalities $\psi_i\le0$ into a bounded zero set of a
quadratic map (bounded because $u$ lies in the ball and each squared
slack equals $-\psi_i(u)$, which is bounded there), which connects the lemma to these results. We use
instead a direct convex KKT elimination and the explicit two-block limit
formula \eqref{eq:limitformula}. The affine-face reduction is what allows
arbitrarily many affine rows. The value bound itself follows from
established algebraic tools; the new part is its application to exact
penalties, with the explicit dependence on the number of nonlinear rows.

\subsection{Separation and multiplier bounds}
\begin{proof}[Proof of Theorem~\ref{thm:fixedcount}]
We bound the separation of every equality-infeasible slice from below and
a linking multiplier of every equality-feasible slice from above, each by
applying Lemma~\ref{lem:nonzerovalue} to an auxiliary problem, and then
apply Lemma~\ref{lem:slices}. Use the objective bound $|f|\le F_0$ from
Section~\ref{sec:upper}. Every auxiliary problem below is obtained by
substituting a boxed integer assignment $z$. Its coefficients keep
bitsize $O(N)$, as in \eqref{eq:height}, so its explicit input length is
at most $N^{c}$ for an absolute constant $c$, uniformly in $z$. Hence the bound $2^{-(N^{c})^{O(k+1)}}$ of
Lemma~\ref{lem:nonzerovalue} is again $2^{-N^{O(k+1)}}$.

\emph{Equality-infeasible slices.}
Let $X_z$ be a nonempty equality-infeasible slice. Let $R_0$ be a rational
bound on all residual magnitudes $|r_j(x,z)|$ over the input box, with bit
length polynomial in $N$. Minimize $t$ over the points $(x,t)$ with
$x\in X_z$ and
\[
 -t\le r_j(x,z)\le t\quad(1\le j\le m),\qquad 0\le t\le R_0.
\]
This is a nonempty compact convex quadratic problem whose only nonlinear
rows are the at most $k$ native rows that are nonlinear in $x$. Its
optimal value is $\delta_z=\min_{X_z}\norm r_\infty$, which is positive
by compactness. Lemma~\ref{lem:nonzerovalue} gives
\begin{equation}\label{eq:fixedseparation}
 \delta_z\ge2^{-N^{O(k+1)}}.
\end{equation}

\emph{Equality-feasible slices: the Slater margin.}
Let $X_z$ be an equality-feasible slice with at least one row nonlinear
in $x$. Let $P_z$ be the polyhedron defined by the native constraints that
are affine in $x$ (including the bounds) and by the relaxed equalities
$r(x,z)=0$. Define the common Slater margin of the nonlinear rows,
\begin{equation}\label{eq:slatermargin}
 \gamma_z=\max\{t:x\in P_z,\ g_i(x,z)+t\le0
           \text{ for all rows $i$ nonlinear in $x$},\ 0\le t\le1\}.
\end{equation}
The refined Slater condition provides a point of $P_z$ at which every
nonlinear row is strict, so $\gamma_z>0$. Minimizing $-t$ in this compact
convex problem with at most $k$ nonlinear rows and applying
Lemma~\ref{lem:nonzerovalue} gives
\begin{equation}\label{eq:marginbound}
 \gamma_z\ge2^{-N^{O(k+1)}}.
\end{equation}

\emph{Multipliers of the nonlinear rows.}
Let $(\bar x,\gamma_z)$ be optimal in \eqref{eq:slatermargin}, and let $x^*$ be an optimum of
the slice problem $v_z=\min\{f(x,z):x\in X_z,\ r(x,z)=0\}$. By the refined
Slater condition, $x^*$ has a KKT certificate, as in
Section~\ref{sec:upper}, with multipliers $\mu_i\ge0$ for
the nonlinear rows, nonnegative multipliers for the native rows affine in
$x$, and linking multipliers for $r=0$. The corresponding convex
Lagrangian is minimized at $x^*$, where its value is $v_z$. Evaluate it
at $\bar x$: the linking terms vanish because $r(\bar x,z)=0$, the terms
of the affine native rows are nonpositive because $\bar x\in P_z$, and
each nonlinear row contributes at most $-\gamma_z\mu_i$. Hence
\begin{equation}\label{eq:nonlinearmult}
 v_z\le f(\bar x,z)-\gamma_z\sum_i\mu_i,\qquad
 \sum_i\mu_i\le\frac{2F_0}{\gamma_z}.
\end{equation}
This is the usual bound on multipliers from a Slater point with a margin.

\emph{Linking multipliers by Cramer's rule.}
It remains to bound the linking multipliers; we do not need to bound the
multipliers of the affine native rows in this particular certificate.
The vector
\[
 b_*=-\nabla_xf(x^*,z)-\sum_i\mu_i\nabla_xg_i(x^*,z)
\]
has coordinates bounded by $2^{N^{O(1)}}(1+2F_0/\gamma_z)$, since all
gradients are affine with coefficients of polynomial bitsize and $x^*$
and $z$ lie in their boxes. By \eqref{eq:marginbound}, this bound is
$2^{N^{O(k+1)}}$. By stationarity, $b_*$ lies in the cone generated by
the gradients of the active native rows affine in $x$ and by both signs
of the linking normals $A_j^T$. By the conic Carath\'eodory theorem,
$b_*$ is a nonnegative combination of linearly independent generators
from this set; let $\nu\le n$ be their number. Their $n\times\nu$
column matrix has rank $\nu$, so some $\nu$ of its rows form a nonsingular
square submatrix. Its entries are rational input coefficients after the integer
substitution, so determinant bounds and Cramer's rule bound the entries
of its inverse by $2^{N^{O(1)}}$. Applying the inverse to the
corresponding $\nu$ coordinates of $b_*$ bounds all selected coefficients
by $2^{N^{O(k+1)}}$. This bound holds for real coefficients; the primal
point $x^*$ and the multipliers $\mu_i$ need not be rational.

Combining the coefficients of oppositely signed linking normals gives a
new KKT certificate at $x^*$: it keeps the multipliers $\mu_i$, puts the
selected coefficients on the active affine rows, and has a linking
multiplier $\lambda_z$ with
\begin{equation}\label{eq:fixedmult}
 \norm{\lambda_z}_1\le2^{N^{O(k+1)}}.
\end{equation}
Only active affine rows received multipliers and the multipliers $\mu_i$
were kept, so stationarity and complementarity still hold. As in
Section~\ref{sec:upper}, the new Lagrangian is convex and minimized at
$x^*$ with value $v_z$; dropping its nonpositive native terms on $X_z$
gives $f(x,z)+\lambda_z^Tr(x,z)\ge v_z$ for all $x\in X_z$. If the slice
has no row nonlinear in $x$, the same argument applies directly with
$b_*=-\nabla_xf(x^*,z)$; KKT conditions hold without a constraint
qualification for affine constraints, and no margin problem is needed.

\emph{Assembly.}
All constants are uniform over the integer assignments. Choose an integer
$\rho$ strictly larger than the bound in \eqref{eq:fixedmult} and than
$2F_0$ times the reciprocal of the bound in \eqref{eq:fixedseparation}.
Its bit length is $N^{O(k+1)}$. The dual norm of $\norm\cdot_\infty$ is
$\norm\cdot_1$, so Lemma~\ref{lem:slices}, with objective bounds
$\underline f=-F_0$ and $\overline f=F_0$, shows that $\rho$ is
minimizer-exact at multiplier zero for the infinity norm. The same $\rho$
works for the one-norm, because $\norm r_1\ge\norm r_\infty$ and both
norms vanish on the original feasible set.
\end{proof}

\begin{remark}
By Theorem~\ref{thm:fixedcount}, superpolynomial penalty encoding within
this model needs a growing number of nonlinear inequalities; arbitrarily
many affine rows alone do not produce it. Some growth with $k$ is
unavoidable: the chain of Remark~\ref{rem:seed} has $k=d-1$ and needs
$\Omega(N2^k)$ bits, against the $N^{O(k+1)}$ bits of
Theorem~\ref{thm:fixedcount}.
\end{remark}

\subsection{A conservative coefficient can be printed in polynomial time}
Because the encoding bounds of Theorems~\ref{thm:generalupper}
and~\ref{thm:fixedcount} are uniform, a sufficient coefficient can be
output without finding an optimum, a Slater point or an active face.

\begin{corollary}\label{cor:printing}
For each fixed continuous dimension $n$, and separately for each fixed
$k\ge0$, there is a polynomial-time algorithm with the following property.
On every input that is promised to satisfy Assumption~\ref{ass:model}
(and, in the second case, to have at most $k$ native inequalities
nonlinear in the continuous variables), it outputs in binary an integer
$\rho$ that is minimizer-exact at multiplier zero for both the infinity
norm and the one-norm.
\end{corollary}
\begin{proof}
Choose a universal integer constant $C$ such that the bounds in
Theorems~\ref{thm:generalupper} and~\ref{thm:fixedcount} are at most
\[
 B_{\mathrm{dim}}=(N+1)2^{C(n+1)}\qquad\text{and}\qquad
 B_{\mathrm{cnt}}=N^{C(k+1)},
\]
respectively. In each case some integer that is minimizer-exact at
multiplier zero has at most $B_{\mathrm{dim}}$, respectively
$B_{\mathrm{cnt}}$, bits and is therefore smaller than
$2^{B_{\mathrm{dim}}}$, respectively $2^{B_{\mathrm{cnt}}}$. By the
consequences of \eqref{eq:fixed}, every larger coefficient is
minimizer-exact at multiplier zero as well. The algorithm therefore prints
a one followed by $B_{\mathrm{dim}}$, respectively $B_{\mathrm{cnt}}$,
zeros. For fixed $n$,
respectively fixed $k$, this output and the time to produce it are
polynomial in $N$.
\end{proof}
For native constraints linear in all variables, a special case of
$k=0$, the analogous results are known: polynomial encoding length for a
jointly convex quadratic objective \citep[Theorem~11]{GuAhmedDey2020},
and polynomial-time output for MILPs with integer data
\citep[Theorem~15]{LefebvreSchmidt2025}, whose Table~1 remarks that the
extension to convex MIQPs is easy (Section~\ref{sec:related}). The case
$k=0$ of Corollary~\ref{cor:printing} also covers rows that are affine in
$x$ but contain products with integer variables, such as
$x_1z_1+z_1^2\le3$. Because all variables are bounded, such rows can
also be rewritten exactly as linear rows of polynomial size, through a
binary expansion of $z$ and exact linearization of the resulting
products, after which \citet[Theorem~11]{GuAhmedDey2020} applies when the
objective is jointly convex; this part of the corollary is therefore
mainly a convenience. The
output of Corollary~\ref{cor:printing} is a
conservative and numerically huge coefficient: the algorithm does not
verify the promises, solve the problem or approximate the least penalty,
and the constant $C$ exists by the cited effective bounds but has not
been evaluated.

\section{Hardness of calibrating the least penalty}\label{sec:calibration}
This section shows that a sufficient coefficient can be trivial to write
down while a nearly least one is hard to compute. On a binary box with one
relaxed equality, a linear objective, a known unique optimizer and the
known sufficient coefficient $\rho=1$, no polynomial-time algorithm
approximates the least penalty within any polynomial factor unless
$\mathrm P=\mathrm{NP}$; the same holds for the zero-multiplier threshold.
The hardness needs a growing number of binary variables; fixed dimension is
not covered.

Throughout, $m=1$ and the norm is the absolute value. On $\R$ every norm is
a positive multiple of $|\cdot|$, which only rescales the thresholds
(Remark~\ref{rem:rescaling}).

\subsection{The calibration problem}
\begin{definition}[Calibration]\label{def:calibration}
Fix a family of instances of \eqref{eq:model} with rational data, encoded
in binary as in Section~\ref{sec:geometry}, and a function $\varphi$ from the
positive integers to $[1,\infty)$. Given an instance of encoding length $N$
whose least penalty $\rho_*$ is finite, \emph{sufficient calibration within
factor $\varphi$} asks for a rational $\widehat\rho$ with
\begin{equation}\label{eq:onesided}
 \rho_*\le\widehat\rho\le \varphi(N)\,\rho_*,
\end{equation}
and \emph{two-sided calibration within factor $\varphi$} asks for a rational
$\widehat\rho$ with
\begin{equation}\label{eq:twosided}
 \rho_*/\varphi(N)\le\widehat\rho\le \varphi(N)\,\rho_*.
\end{equation}
Calibration at multiplier zero is defined in the same way, with
$\rho_*(0)$ in place of $\rho_*$. \emph{Polynomial-factor calibration}
means that $\varphi=p$ is a fixed polynomial with $p(N)\ge1$ for all $N$.
\end{definition}

A one-sided estimate \eqref{eq:onesided} is itself a sufficient
coefficient: the value-exact coefficients form the interval
$[\rho_*,\infty)$, and the coefficients exact at multiplier zero form
$[\rho_*(0),\infty)$ (Corollary~\ref{cor:ratio} and the discussion after
it). A two-sided estimate need not be sufficient. On the families below a
sufficient coefficient is known in advance, so the only question is how
close to the least one an algorithm can get.

\subsection{A binary-box family}\label{sec:calibration-family}
Take a \textsc{Subset Sum} instance: positive integers $w_1,\ldots,w_t$ and
a positive integer target $\beta$, with encoding length $N_0$. For
$I\subseteq\{1,\ldots,t\}$ write $w(I)=\sum_{i\in I}w_i$. The instance is a
YES instance if $w(I)=\beta$ for some $I$, and a NO instance otherwise.
Deciding this is NP-complete. Append the item $w_{t+1}=\beta+1$, extend the
notation $w(I)$ to $I\subseteq\{1,\ldots,t+1\}$, choose an integer
$\kappa\ge2$, and set
\begin{equation}\label{eq:subsetmodel}
 X=\{0,1\}^{t+2}\ni(x,q),\qquad f(x,q)=-q,\qquad
 r(x,q)=\kappa\sum_{i=1}^{t+1}w_ix_i-(\kappa\beta+1)q.
\end{equation}
The instance is given by the integer coefficient vector
$(\kappa w_1,\ldots,\kappa w_{t+1},-(\kappa\beta+1))$ of $r$; the box and the
objective are implicit.

\paragraph{How the gadget works.}
Read $x$ as the indicator vector of a subset $I$ and $q$ as a switch
whose slice $q=1$ is improving: it lowers the objective from $0$ to $-1$. The residual is
$\kappa w(I)\ge0$ at $q=0$ and $\kappa(w(I)-\beta)-1$ at $q=1$. By
Proposition~\ref{prop:scalar}, the least penalty is governed by how close
the improving slice $q=1$ comes to $r=0$ from each side. Each part of the
construction has one purpose.
\begin{itemize}
 \item The $+1$ in $\kappa\beta+1$ makes the whole slice $q=1$
 equality-infeasible: $r=0$ at $q=1$ would require $\kappa$ to divide
 $\kappa\beta+1$. At $q=0$, $r=0$ forces $x=0$ because all items are
 positive. Hence $(x,q)=(0,0)$ is the unique feasible point and $v=0$,
 known without solving \textsc{Subset Sum}.
 \item The appended item $\beta+1$ fixes the nearest positive residual at
 $q=1$ to $\kappa-1$ on every instance, so the positive side carries no
 information about the instance.
 \item The factor $\kappa$ separates the cases on the negative side: a
 subset with $w(I)=\beta$ gives residual $-1$, while every subset with
 $w(I)\le\beta-1$ gives residual at most $-(\kappa+1)$.
\end{itemize}
Linear optimization over $X$ is trivial; this does not make the penalized
subproblem with the absolute residual easy.

Formally, let $\beta_-=\max\{w(I):I\subseteq\{1,\ldots,t+1\},\
w(I)\le\beta\}$ be the largest attainable sum not exceeding the target
(the empty set shows that it exists), and define
\begin{equation}\label{eq:nearestsums}
 g_-=\kappa(\beta-\beta_-)+1,\qquad g_+=\kappa-1.
\end{equation}
At $q=1$ the residual $\kappa(w(I)-\beta)-1$ is negative exactly when
$w(I)\le\beta$; its magnitude is then $\kappa(\beta-w(I))+1\ge g_-$, with
equality when $w(I)=\beta_-$. The residual is positive exactly when
$w(I)\ge\beta+1$; it is then at least $\kappa-1=g_+$, with equality for
$I=\{t+1\}$. No subset with sum at most $\beta$ contains the item $t+1$, so
$\beta_-$ is the best sum of the original instance. Consequently
\begin{equation}\label{eq:gapcases}
 g_-=1\ \text{in YES instances},\qquad g_-\ge\kappa+1\ \text{in NO instances}.
\end{equation}

\begin{proposition}\label{prop:calibrationdual}
Let $w_1,\ldots,w_t,\beta$ be positive integers, $w_{t+1}=\beta+1$, and
$\kappa\ge2$ an integer. The instance \eqref{eq:subsetmodel} has $v=0$,
attained only at $(x,q)=(0,0)$. For every $\rho\ge0$,
\begin{equation}\label{eq:calibrationdual}
 D_\rho=\min\left\{0,-1+\rho\,\frac{2g_-g_+}{g_-+g_+}\right\},
\end{equation}
and the multiplier $\lambda_\rho=\rho(g_--g_+)/(g_-+g_+)$ attains this
value. Consequently
\begin{equation}\label{eq:calibrationthresholds}
 \rho_*=\frac12\left(\frac1{g_-}+\frac1{g_+}\right),\qquad
 \rho_*(0)=\max\left\{\frac1{g_-},\frac1{g_+}\right\}.
\end{equation}
\end{proposition}
\begin{proof}
The plan is to bound $D_\rho$ above by a two-point balanced mixture and to
show that $\lambda_\rho$ attains this bound. The value $v=0$ and the
uniqueness of the feasible point were shown above.

Let $P_-$ be a point with $q=1$ and $w(I)=\beta_-$, and let $P_+$ be the
point with $q=1$ and $I=\{t+1\}$. Then $(f,r)=(-1,-g_-)$ at $P_-$ and
$(f,r)=(-1,g_+)$ at $P_+$. The weights $g_+/(g_-+g_+)$ on $P_-$ and
$g_-/(g_-+g_+)$ on $P_+$ give a balanced mixture: its average residual is
zero. Its average of $f+\rho|r|$ is $-1+2\rho g_-g_+/(g_-+g_+)$, so
Proposition~\ref{prop:mixture} bounds $D_\rho$ by this number. Weak duality
bounds $D_\rho$ by $v=0$. This proves ``$\le$'' in
\eqref{eq:calibrationdual}.

For the reverse inequality, note that
$\rho-\lambda_\rho=2\rho g_+/(g_-+g_+)\ge0$ and
$\rho+\lambda_\rho=2\rho g_-/(g_-+g_+)\ge0$. A native point with $q=1$ and
$r<0$ has $|r|\ge g_-$, so its penalized value at $\lambda_\rho$ is
\[
 -1+(\rho-\lambda_\rho)|r|\ge-1+(\rho-\lambda_\rho)g_-
 =-1+\rho\,\frac{2g_-g_+}{g_-+g_+}.
\]
A native point with $q=1$ and $r>0$ has $r\ge g_+$, and
$-1+(\rho+\lambda_\rho)r\ge-1+(\rho+\lambda_\rho)g_+$ is the same number.
A native point with $q=0$ has $f=0$ and $r\ge0$, so its value
$(\rho+\lambda_\rho)r$ is nonnegative. Hence $L_\rho(\lambda_\rho)$ is at
least the right-hand side of \eqref{eq:calibrationdual}, and the upper
bound is attained.

By \eqref{eq:calibrationdual}, $D_\rho=0$ exactly when
$\rho\ge(g_-+g_+)/(2g_-g_+)$, which is the formula for $\rho_*$. For
$\rho_*(0)$, apply \eqref{eq:fixed} with $\lambda=0$ and $v=0$: an
infeasible point with $q=0$ contributes $0/|r|=0$, and an infeasible point
with $q=1$ contributes $1/|r|$, which is largest at the smallest residual
magnitudes $g_-$ and $g_+$.
\end{proof}

The formulas agree with Proposition~\ref{prop:scalar}, where
$A_+=1/g_+$ (points with $q=0$ contribute $0$) and $A_-=1/g_-$. At
$\rho=\rho_*$ the interval of exact multipliers in \eqref{eq:scalar}
reduces to the single multiplier
$\lambda_{\rho_*}=(1/g_+-1/g_-)/2$, and both witnesses $P_-$ and $P_+$ tie
with the optimum. So $\rho_*$ is exact at $\lambda_{\rho_*}$ but not
minimizer-exact. Every $\rho>\rho_*$ is minimizer-exact at
$\lambda_{\rho_*}$: then $\rho-\lambda_{\rho_*}>1/g_-$ and
$\rho+\lambda_{\rho_*}>1/g_+$, so every infeasible point has a positive
penalized value. On this family $1/(2\kappa)<\rho_*\le1$, because
$g_+=\kappa-1$ and $g_-,g_+\ge1$.

\paragraph{A sufficient coefficient is known.}
For every instance of the family, $\rho=1$ is exact at multiplier zero.
Every nonzero residual is an integer and the objective is $0$ or $-1$, so
$f+|r|\ge0$ at every native point. An infeasible point has $f+|r|=0$ only
if $q=1$ and $|r|=1$. The residual $-1$ occurs exactly in YES instances,
and the residual $+1$ occurs exactly when $\kappa=2$ (through
$I=\{t+1\}$). These ties are the only obstruction: every $\rho>1$ is
minimizer-exact at multiplier zero.

\paragraph{The YES/NO gap.}
By \eqref{eq:gapcases} and \eqref{eq:calibrationthresholds},
\begin{equation}\label{eq:thresholdgap}
 \begin{aligned}
 &\rho_*=\frac{\kappa}{2(\kappa-1)}>\frac12,\qquad
 \rho_*(0)=1 &&\text{in YES instances},\\
 &0<\rho_*\le\frac{\kappa}{\kappa^2-1},\qquad
 \rho_*(0)=\frac1{\kappa-1} &&\text{in NO instances}.
 \end{aligned}
\end{equation}
The ratio between the YES value and the NO bound is $(\kappa+1)/2$ for
$\rho_*$ and $\kappa-1$ for $\rho_*(0)$. For $\kappa=2$ the zero-multiplier
threshold equals $1$ in both cases, so part~(b) below uses $\kappa\ge3$.

\subsection{Testing exactness and calibrating are hard}
The first result shows that testing a single given coefficient is already
hard. The membership argument uses that $v=0$ is known on this family; for
a general instance, certifying $D_\rho<v$ would also require
information about $v$.

\begin{theorem}\label{thm:coNP}
Consider the instances \eqref{eq:subsetmodel} with $\kappa=4$, each given
by the integer coefficient vector $(4w_1,\ldots,4w_{t+1},-(4\beta+1))$ of
$r$. Deciding whether $\rho=1/3$ is value-exact, that is, whether
$D_{1/3}=0$, is coNP-complete. More precisely, $D_{1/3}=-1/2$ for YES
instances of \textup{\textsc{Subset Sum}} and $D_{1/3}=0$ for NO instances.
\end{theorem}
\begin{proof}
For YES instances, $g_-=1$ and $g_+=3$, and \eqref{eq:calibrationdual} gives
$D_{1/3}=\min\{0,-1+\tfrac13\cdot\tfrac64\}=-1/2$. For NO instances,
$\rho_*\le4/15<1/3$ by \eqref{eq:thresholdgap}, so $D_{1/3}=0$.

Hardness: the map from a \textsc{Subset Sum} instance to the coefficient
vector is computable in polynomial time and sends NO instances to exact
coefficients and YES instances to non-exact ones. It is therefore a
polynomial reduction from the complement of \textsc{Subset Sum}, which is
coNP-complete.

Membership in coNP: whether an input vector belongs to the family can be
checked in polynomial time, since $\beta$, the items $w_i$ and the
condition $w_{t+1}=\beta+1$ are recovered by integer division by $4$. On a
member of the family, $D_{1/3}<0$ holds exactly in YES instances, and a
subset $I\subseteq\{1,\ldots,t\}$ with $w(I)=\beta$ is a polynomial
certificate of this.
\end{proof}
Consequently, computing $\rho_*$ on this family is NP-hard, and evaluating
$D_{1/3}$ within absolute error strictly below $1/4$ decides
\textsc{Subset Sum}. In contrast, $D_0=-1$ on every instance: without the
norm term the Lagrangian dual is trivial to evaluate and never closes the
gap.

The second result says that even approximating the threshold, from above,
within a polynomial factor is hard.

\begin{theorem}\label{thm:polyfactor}
Let $p$ be a fixed polynomial with $p(N)\ge1$ for all $N$. Unless
$\mathrm P=\mathrm{NP}$, no polynomial-time algorithm achieves sufficient
calibration within factor $p$ on the family \eqref{eq:subsetmodel}, either
\begin{enumerate}
 \item[(a)] for the least penalty: on every instance of encoding length
 $N$, return a rational $\widehat\rho$ with
 \begin{equation}\label{eq:polyfactor}
  \rho_*\le\widehat\rho\le p(N)\,\rho_*;
 \end{equation}
 \item[(b)] or at multiplier zero: on every instance of encoding length $N$,
 return a rational $\widehat\rho$ with
 $\rho_*(0)\le\widehat\rho\le p(N)\,\rho_*(0)$, where
 $\rho_*(0)=\max\{1/g_-,1/g_+\}$; this equals $1$ in YES instances and
 $1/(\kappa-1)$ in NO instances.
\end{enumerate}
\end{theorem}
\begin{proof}
The plan is to choose $\kappa$ exponentially large, so that the YES/NO
ratio in \eqref{eq:thresholdgap} exceeds $p(N)$, and then to read off the
\textsc{Subset Sum} answer by comparing $\widehat\rho$ with a fixed number.

Given a \textsc{Subset Sum} instance of length $N_0$, take
$\kappa=2^{N_0}$. The constructed instance has $t+2\le N_0+2$ coefficients,
each with $O(N_0)$ bits, so its length is $N=O(N_0^2)$ and it is computed
in polynomial time. In particular $p(N)$ is bounded by a polynomial in
$N_0$.

(a) In YES instances every admissible output satisfies
$\widehat\rho\ge\rho_*>1/2$. In NO instances every admissible output
satisfies
\[
 \widehat\rho\le p(N)\,\frac{\kappa}{\kappa^2-1}
 =p(N)\,\frac{2^{N_0}}{4^{N_0}-1}<\frac12
\]
for all sufficiently large $N_0$, since the exponential term dominates.
(b) For $N_0\ge2$ we have $\kappa\ge3$. In YES instances every admissible
output satisfies $\widehat\rho\ge1$; in NO instances
$\widehat\rho\le p(N)/(2^{N_0}-1)<1$ for all sufficiently large $N_0$.

In both cases there is a constant $n_1$, depending only on $p$, such that
the comparison of $\widehat\rho$ with $1/2$ (respectively $1$) decides
every \textsc{Subset Sum} instance with $N_0\ge n_1$. Instances with
$N_0<n_1$ are finitely many and are decided by exhaustive enumeration. The
output $\widehat\rho$ of a polynomial-time algorithm has polynomially many
bits, so the comparison takes polynomial time. This would decide
\textsc{Subset Sum} in polynomial time.
\end{proof}

\begin{remark}[Two-sided estimates and constant factors]\label{rem:twosided}
Two-sided calibration is not easier. If $\widehat\rho$ satisfies
\eqref{eq:twosided} with $\varphi=p$, replace $p$ by a pointwise larger polynomial
$p'$ with positive integer coefficients (so that $p'(N)$ is an integer
computable in polynomial time); then $p'(N)\widehat\rho$ satisfies
\eqref{eq:onesided} with $\varphi=p'^2$, again a polynomial. Thus
Theorem~\ref{thm:polyfactor} also rules out two-sided polynomial-factor
calibration, in both parts. For a fixed constant factor $C\ge1$, an
exponential $\kappa$ is not needed. In part~(a), a one-sided estimate within
factor $C$ separates the YES value $\kappa/(2(\kappa-1))$ from the NO bound
$C\kappa/(\kappa^2-1)$ exactly when $C<(\kappa+1)/2$, so any constant
$\kappa$ with $\kappa+1>2C$ suffices. In part~(b), any constant $\kappa$ with
$\kappa-1>C$ suffices.
\end{remark}

\begin{remark}[Subexponential factors]\label{rem:subexpfactor}
The proof gives more than polynomial factors. For every fixed
$\epsilon\in(0,1)$, sufficient calibration within factor
$\varphi(N)=2^{N^{1-\epsilon}}$ is also impossible in polynomial time unless
$\mathrm P=\mathrm{NP}$, in both parts. Choose an integer
$c>(1-\epsilon)/\epsilon$ and take $\kappa=2^{N_0^c}$. Each coefficient
then has at most $N_0^c+O(N_0)$ bits, so $N=O(N_0^{c+1})$ and the
construction is polynomial. Since $(c+1)(1-\epsilon)<c$, we have
$N^{1-\epsilon}\le N_0^c-2$, that is $\varphi(N)\le\kappa/4$, for all
sufficiently large $N_0$. Then in part~(a) every NO output satisfies
$\widehat\rho\le \varphi(N)\kappa/(\kappa^2-1)\le\kappa^2/(4(\kappa^2-1))<1/2$,
while every YES output exceeds $1/2$. In part~(b) every NO output satisfies
$\widehat\rho\le \varphi(N)/(\kappa-1)\le\kappa/(4(\kappa-1))<1$, while every YES
output is at least $1$.
\end{remark}

\begin{remark}[Rescaling does not help]\label{rem:rescaling}
Replacing $r$ by $cr$ for a rational $c>0$ replaces $L_\rho(\lambda)$ by
$L_{c\rho}(c\lambda)$ of the original instance; replacing the absolute value
by the norm $c|\cdot|$ replaces it by $L_{c\rho}(\lambda)$. In both cases $\rho_*$ and $\rho_*(0)$ are divided by $c$, and the
YES/NO ratios $(\kappa+1)/2$ and $\kappa-1$ are unchanged. A guarantee
within a factor is invariant under such rescaling, so no choice of scaling
makes calibration easier on this family.
\end{remark}

\begin{remark}[Weak NP-hardness]\label{rem:weak}
The reduction is from \textsc{Subset Sum}, which is NP-hard only when its
numbers are large, and the constructed coefficients are large as well. This
is unavoidable on this family: a dynamic program over the attainable values
of $\sum_iw_ix_i$ computes $g_-$, hence $\rho_*$ and $\rho_*(0)$ exactly, in
time polynomial in $N$ and the largest coefficient magnitude. So the family
gives weak NP-hardness only. Proposition~\ref{prop:stableset} below gives
strong NP-hardness of exact computation with coefficients in
$\{-1,0,1\}$, at the price of a native set over which linear optimization
is itself hard.
\end{remark}

\begin{remark}[Where the hardness lies]\label{rem:hardnesslocation}
The difficulty is testing whether a given coefficient is exact
(Theorem~\ref{thm:coNP}), not locating the threshold once such a test is
available. Suppose an oracle decides, for a given instance of the family and
a rational $\rho$, whether $\rho$ is value-exact. Since
$\rho_*\in(1/(2\kappa),1]$ and the value-exact coefficients form
$[\rho_*,\infty)$, testing $\rho=2^{-j}$ for
$j=0,1,\ldots,\lceil\log_2\kappa\rceil+1$ finds the largest $j$ with
$2^{-j}$ exact; then $2^{-j}\ge\rho_*>2^{-j-1}$, so $\widehat\rho=2^{-j}$ is
a sufficient calibration within factor two, after $O(\log\kappa)=O(N)$
oracle calls. The same works at multiplier zero with an oracle for
$L_\rho(0)=v$. In decomposition methods the penalized subproblem is solved
anyway. The results above say that nearly least penalties are as hard to
certify as those subproblems, even when the original problem is trivial.
\end{remark}

\subsection{Unit coefficients with a hard native set}
The binary-box family has large numbers. The next family has all
coefficients in $\{-1,0,1\}$, and its least penalty is the stability number
of a graph. This gives strong NP-hardness of exact computation, because the
native set now encodes a hard optimization problem.

For a simple graph $G=(V,E)$ with $V\ne\varnothing$, let $\alpha(G)$ be the
maximum size of a stable set. Take binary variables $x_i$ ($i\in V$),
$y_+$ and $y_-$, with native constraints
\begin{equation}\label{eq:stablemodel}
 y_++y_-\le1,\qquad x_i\le y_++y_-\ (i\in V),\qquad
 x_i+x_j\le1\ (ij\in E),
\end{equation}
objective $f=-\sum_{i\in V}x_i$ and relaxed equality $r=y_+-y_-=0$.

\begin{proposition}\label{prop:stableset}
For the instance \eqref{eq:stablemodel}, $v=0$ is attained only at the
zero vector, and for all $\rho\ge0$ and $\lambda\in\R$,
\begin{equation}\label{eq:stabledual}
 L_\rho(\lambda)=\min\{0,\rho-\alpha(G)-|\lambda|\},\qquad
 D_\rho=\min\{0,\rho-\alpha(G)\},\qquad \rho_*=\rho_*(0)=\alpha(G).
\end{equation}
At $\rho=\alpha(G)$, $\lambda=0$ is the only exact multiplier and the
infeasible points given by maximum stable sets tie with the optimum. At
multiplier zero, $\rho$ is minimizer-exact if and only if $\rho>\alpha(G)$.
Consequently, computing $\rho_*$ is strongly NP-hard. Given $G$ and an
integer $c$ as input, deciding whether $\rho=c$ is value-exact, that is,
whether $\alpha(G)\le c$, is coNP-complete; a stable set of size $c+1$
certifies failure. For each fixed $c$, enumerating the $(c+1)$-subsets of
$V$ decides it in polynomial time.
\end{proposition}
\begin{proof}
If $r=0$, then $y_+=y_-$, and $y_++y_-\le1$ forces $y_+=y_-=0$; then
$x_i\le0$ for all $i\in V$. So the zero vector is the unique feasible point and
$v=0$. The native points form three slices. If $y_+=y_-=0$, then $x=0$
and $(f,r)=(0,0)$. If $(y_+,y_-)=(1,0)$, the constraints on $x$ say exactly
that $x$ is the indicator vector of a stable set $S$, and
$(f,r)=(-|S|,1)$. If $(y_+,y_-)=(0,1)$, likewise $(f,r)=(-|S|,-1)$.
Minimizing $f+\lambda r+\rho|r|$ over the three slices gives
$\min\{0,\rho-\alpha(G)+\lambda,\rho-\alpha(G)-\lambda\}$, which is the
formula for $L_\rho(\lambda)$. Its supremum over $\lambda$ is attained at
$\lambda=0$, which gives $D_\rho$, and $D_\rho=0$ exactly when
$\rho\ge\alpha(G)$; the same holds for $L_\rho(0)$, so
$\rho_*=\rho_*(0)=\alpha(G)$.

At $\rho=\alpha(G)$, $L_\rho(\lambda)=\min\{0,-|\lambda|\}$ equals $0$ only
for $\lambda=0$, and a maximum stable set with $y_+=1$ then has penalized
value $0$. At multiplier zero, every infeasible point has value
$\rho-|S|$ for a stable set $S$, which is positive for all such points
exactly when $\rho>\alpha(G)$.

All coefficients and right-hand sides lie in $\{-1,0,1\}$, and computing
$\alpha(G)$ is NP-hard, so computing $\rho_*$ is strongly NP-hard.
Deciding $\alpha(G)\le c$ for an input integer $c$ is the complement of the
NP-complete stable-set decision problem, and a stable set of size $c+1$ is
a polynomial certificate that $\rho=c$ is not value-exact. For fixed $c$
there are $O(|V|^{c+1})$ subsets of size $c+1$ to check.
\end{proof}
Unlike the binary box, where linear optimization over $X$ is trivial, here
linear optimization over the native set is already the maximum stable-set
problem. We use this family only for strong NP-hardness of exact
computation. It cannot exclude every polynomial factor as
Theorem~\ref{thm:polyfactor} does: since $1\le\alpha(G)\le|V|\le N$, the output $\widehat\rho=|V|$ is
always a sufficient calibration within factor $|V|\le N$. The numeric
family \eqref{eq:subsetmodel} is what yields arbitrary polynomial factors.

\subsection{Relation to prior work}
In their conclusion, \citet{LefebvreSchmidt2025} ask whether the smallest
penalty parameter that closes the duality gap can be computed in
polynomial time, and conjecture a negative answer. Their gap is defined
through the supremum over multipliers (their Definition~1), which is
value-exactness here, and their question concerns MILPs and MIQPs. The
binary-box family \eqref{eq:subsetmodel} is a pure binary integer linear
program, so it lies within the scope of their question.
Theorem~\ref{thm:polyfactor}(a) confirms their conjecture unless
$\mathrm P=\mathrm{NP}$, even for approximation within any polynomial
factor, in growing binary dimension. Section~\ref{sec:related} compares
the results with earlier hardness results for choosing penalties.

\section{Right-hand-side perturbations}\label{sec:perturbation}
The coefficient in Theorem~\ref{thm:lower} explodes because the
right-hand side of the relaxed equality lies within $\delta_d$ of the
boundary of the residual image of an improving slice, that is, of the
set of right-hand sides attainable on that slice. This section shows
that a margin from such boundaries controls the coefficient: if the
perturbed problem is feasible and its right-hand side keeps distance more
than $\gamma$ from the boundaries of all residual images, the coefficient
$M/\gamma$ is exact at multiplier zero (Lemma~\ref{lem:rhsmargin}). A right-hand side sampled from a fine
rational grid has such a margin with high probability, which gives a
coefficient of polynomial encoding length
(Theorem~\ref{thm:gridsmoothing}). The principle is classical in
conditioning theory \citep{DunaganSpielmanTeng2011,BuergisserAmelunxen2012}:
a right-hand side far from the boundaries of the convex residual images
is far from any change in the set of feasible integer assignments, and
random perturbations rarely come close to such boundaries. We prove the
cube estimate and its transfer to the grid; the Gaussian variant in
Section~\ref{sec:gaussian} uses \citet[Lemma~2.3.4]{DunaganSpielmanTeng2011}.

Perturbing the right-hand side changes the problem. All guarantees in
this section concern the sampled instance, which may be infeasible and
may have a different optimal integer assignment from the original one.

\subsection{Setting}\label{sec:perturbation-setting}
Let $\mathcal Z$ be a finite index set with $|\mathcal Z|\le K$, where
$K\ge1$. For each $z\in\mathcal Z$, let $X_z$ be a nonempty compact
convex subset of a Euclidean space (the slice), let
$r_z:X_z\to\R^m$ be affine, with $m\ge1$, and let $f_z:X_z\to\R$ be
continuous and convex. Bounds $\underline f\le\overline f$ with
\[
 \underline f\le f_z(x)\le\overline f\qquad(z\in\mathcal Z,\ x\in X_z)
\]
are supplied, and $M=\overline f-\underline f$ denotes the objective range.
The slices, maps and objectives are fixed before any sampling. For a
right-hand side $b\in\R^m$, let
\begin{equation}\label{eq:rhsmodel}
 \begin{aligned}
 v(b)&=\min\{f_z(x):z\in\mathcal Z,\ x\in X_z,\ r_z(x)=b\},\\
 p_\rho(b)&=\min_{z\in\mathcal Z,\,x\in X_z}
             \{f_z(x)+\rho\norm{r_z(x)-b}_\infty\},
 \end{aligned}
\end{equation}
with $v(b)=+\infty$ if no point is feasible. We call $b$ \emph{feasible}
if $v(b)<\infty$, and we write $v_z(b)$ for the minimum of $f_z$ over
$\{x\in X_z:r_z(x)=b\}$ (with $v_z(b)=+\infty$ if this set is empty).

This is the setting of Section~\ref{sec:geometry}: the native set is the
disjoint union of the slices, the objective is $f_z$ on slice $z$, the
residual is $r_z(x)-b$, and the norm is the infinity norm. In particular,
$p_\rho(b)=L_\rho(0)$ for this instance. For a feasible $b$, we write
$\rho_*[b]$ for the least penalty of this instance, with the multiplier
optimized; the square brackets distinguish it from the fixed-multiplier
threshold $\rho_*(\lambda)$. For infeasible $b$ the least penalty is not
defined; where the distribution of $\rho_*[b]$ over random $b$ is
discussed, any nonnegative value may be assigned to infeasible $b$.

For example, in the model of Section~\ref{sec:upper} with convex slices,
each integer assignment $z$ with a nonempty slice gives such a slice,
with $f_z$ the objective and $r_z$ the left-hand side of the linking
equality for this fixed $z$. Then $K$ can be taken as the product of the
numbers of allowed values of the integer variables, so $\log K$ is
polynomial in the input length, but $K$ itself can be exponential in the
number of integer variables. The present setting needs neither a Slater condition nor
quadratic data.

The \emph{residual image} $C_z=r_z(X_z)$ is the set of right-hand sides
attainable on slice $z$; it is nonempty, compact and convex, and the
problem with right-hand side $b$ has a feasible point in slice $z$
exactly when $b\in C_z$.
Boundaries $\partial C$ are topological boundaries in $\R^m$. Thus a
closed convex set with empty interior, such as a singleton, is its own
boundary, and every right-hand side in it has margin zero: arbitrarily
small perturbations leave it. Distances are
$\operatorname{dist}_\infty(b,S)=\inf_{y\in S}\norm{b-y}_\infty$, and
similarly $\operatorname{dist}_2$, with distance $+\infty$ to the empty
set. The latter convention matters only in Lemma~\ref{lem:cubetube},
where $C=\varnothing$ or $C=\R^m$ is allowed.

\subsection{A deterministic boundary-margin bound}
The idea is a convex repair. Suppose every residual image has boundary
distance greater than $\gamma$ from $b$. On a slice whose image contains
$b$, the whole cube $b+[-\gamma,\gamma]^m$ then consists of attainable
right-hand sides; in particular, this image is full-dimensional. A native
point $x$ at residual distance $t=\norm{r_z(x)-b}_\infty>0$ can be
repaired by mixing it with a point $x'$ whose residual lies on the far
side of $b$ at distance $\gamma$. The mixture has residual exactly $b$,
and convexity bounds the cost of the repair by $(M/\gamma)t$. On a slice
whose image does not contain $b$, every residual is at distance more than
$\gamma$ from $b$, so the penalty $(M/\gamma)t$ exceeds the whole
objective range.

\begin{lemma}\label{lem:rhsmargin}
In the setting of Section~\ref{sec:perturbation-setting}, let $b\in\R^m$
be feasible and let $\gamma>0$ satisfy
$\operatorname{dist}_\infty(b,\partial C_z)>\gamma$ for every
$z\in\mathcal Z$. Then $p_{M/\gamma}(b)=v(b)$, that is, the coefficient
$M/\gamma$ is exact at multiplier zero for right-hand side $b$; in
particular $\rho_*[b]\le M/\gamma$. Every coefficient $\rho>M/\gamma$ is
minimizer-exact at multiplier zero.
\end{lemma}
\begin{proof}
We show that $f_z(x)+(M/\gamma)\norm{r_z(x)-b}_\infty\ge v(b)$ for every
slice $z$ and every $x\in X_z$, first on slices whose image contains $b$
and then on the others.

\emph{Slices with $b\in C_z$.} First, $b+[-\gamma,\gamma]^m\subseteq C_z$.
Otherwise some $y$ with $\norm{y-b}_\infty\le\gamma$ lies outside $C_z$.
Since $C_z$ is closed and contains $b$, the last point of $C_z$ on the
segment from $b$ to $y$ lies on $\partial C_z$, at infinity-norm distance
at most $\gamma$ from $b$, a contradiction.

Now let $x\in X_z$ with $t=\norm{r_z(x)-b}_\infty>0$. The point
$b-(\gamma/t)(r_z(x)-b)$ lies in $b+[-\gamma,\gamma]^m\subseteq C_z$, so
some $x'\in X_z$ satisfies
\[
 r_z(x')=b-\frac\gamma t\,(r_z(x)-b).
\]
The convex combination $x_b=(\gamma x+tx')/(\gamma+t)$ lies in $X_z$, and
since $r_z$ is affine, $r_z(x_b)=b$. Convexity of $f_z$ and
$f_z(x')\le\overline f$ give
\[
 v_z(b)\le f_z(x_b)\le\frac\gamma{\gamma+t}f_z(x)
                    +\frac t{\gamma+t}\overline f .
\]
Rearranging, and using $v(b)\le v_z(b)$ and
$\overline f-v_z(b)\le\overline f-\underline f=M$,
\[
 f_z(x)\ge v_z(b)-\frac{\overline f-v_z(b)}{\gamma}\,t
       \ge v(b)-\frac M\gamma\,t .
\]
This inequality also holds at $t=0$, since then $x$ is feasible and
$f_z(x)\ge v(b)$.

\emph{Slices with $b\notin C_z$.} A closest point of the compact set
$C_z$ to $b$ lies on $\partial C_z$, so
$\operatorname{dist}_\infty(b,C_z)=\operatorname{dist}_\infty(b,\partial C_z)>\gamma$.
Hence every $x\in X_z$ has $t=\norm{r_z(x)-b}_\infty>\gamma$ and
\[
 f_z(x)+\frac M\gamma\,t\ge\underline f+M=\overline f\ge v(b).
\]

Thus $p_{M/\gamma}(b)\ge v(b)$, and a feasible optimum attains $v(b)$
with zero penalty, so equality holds. If $\rho>M/\gamma$, every point
with nonzero residual has penalized value strictly greater than $v(b)$,
so the minimizers are exactly the optimal feasible points.
\end{proof}
The argument needs neither differentiability nor computed multipliers.
It is a right-hand-side analogue of Lemma~\ref{lem:slices}: the margin
$\gamma$ replaces the multiplier bound on slices whose image contains
$b$ by the convex repair, and it gives the residual separation $\gamma$
on the other slices. The relation between exact penalties and multiplier bounds in
convex optimization is classical; see
\citet[Theorem~4.2]{FriedlanderTseng2007} for a polar-gauge
characterization under a nonempty compact solution set.

Convexity of the objective is essential. Take one slice $X=[-1,1]$ with
$r(x)=x$, right-hand side $b=0$ and $f(x)=-\sqrt{|x|}$. The image
$[-1,1]$ has boundary distance one from $b$, and the objective range is
one. For $0<t\le1$, the balanced mixture of $x=t$ and $x=-t$ with equal
weights has average penalized value $-\sqrt t+\rho t$, independently of
the multiplier. This is negative for small $t$, so by
Proposition~\ref{prop:mixture}, $D_\rho<0=v$ for every $\rho\ge0$: no
finite penalty is value-exact.

\subsection{Convex-boundary tubes in a cube}
The next lemma bounds the probability that a uniform point of a cube
falls near the boundary of a closed convex set. The bound is linear in
the tube width and does not depend on the set.

\begin{lemma}\label{lem:cubetube}
Let $Q=b_0+[-\sigma,\sigma]^m$ with $\sigma>0$, and let $B$ be uniform on
$Q$. For every closed convex set $C\subseteq\R^m$ and every $\gamma\ge0$,
\begin{equation}\label{eq:cubetube}
 \Pr\{\operatorname{dist}_\infty(B,\partial C)\le\gamma\}
       \le\min\{1,2m\gamma/\sigma\}.
\end{equation}
The same bound holds with $\operatorname{dist}_2$ in place of
$\operatorname{dist}_\infty$.
\end{lemma}
\begin{proof}
The plan is to cover the tube by an outer shell and an inner shell of the
set, bound the volume of each shell inside $Q$ one coordinate direction
at a time, and let the shell width decrease to $\gamma$.

If $C=\varnothing$ or $C=\R^m$, the boundary is empty and the event is
empty. Otherwise $\partial C$ is nonempty. Fix $w>0$ and let
$S_i=[-w,w]e_i$ for $i=1,\ldots,m$. For a closed convex set $E$, the
addition $E+S_i$ and the erosion $E\ominus S_i=\{x:x+S_i\subseteq E\}$
are again closed and convex.

\emph{One-dimensional fact.} Let $I\subseteq\R$ be a closed convex set
(empty, a point, a closed interval, a closed ray or $\R$) and let $W$ be
an interval of length $2\sigma$. Then
\[
 \bigl|((I+[-w,w])\setminus I)\cap W\bigr|\le2w,\qquad
 \bigl|(I\setminus(I\ominus[-w,w]))\cap W\bigr|\le2w,
\]
where $|\cdot|$ is length. Indeed, addition extends $I$ by $w$ at each
finite endpoint, and erosion removes at most $w$ at each finite endpoint
(possibly deleting a short interval entirely); there are at most two
endpoints.

\emph{Shell volumes.} On each line parallel to $e_i$, the section of
$E+S_i$ is the section of $E$ plus $[-w,w]$, the section of $E\ominus S_i$
is the section of $E$ eroded by $[-w,w]$, and the section of $Q$ is an
interval of length $2\sigma$. The one-dimensional fact and Fubini's
theorem over the $(m-1)$-dimensional face of $Q$ orthogonal to $e_i$ give
\begin{align*}
 \operatorname{vol}\bigl(((E+S_i)\setminus E)\cap Q\bigr)
   &\le2w(2\sigma)^{m-1},\\
 \operatorname{vol}\bigl((E\setminus(E\ominus S_i))\cap Q\bigr)
   &\le2w(2\sigma)^{m-1}.
\end{align*}
We measure set differences intersected with $Q$, which have finite volume.
Successive additions of $S_1,\ldots,S_m$ to $C$ give $C+[-w,w]^m$.
Successive erosions give $C\ominus[-w,w]^m$, because
$(E\ominus S)\ominus T=E\ominus(S+T)$. The intermediate sets are closed
and convex, so the bounds apply at each step, and telescoping over the
$m$ steps gives
\[
 \operatorname{vol}\bigl(((C+[-w,w]^m)\setminus C)\cap Q\bigr)\le2mw(2\sigma)^{m-1},
 \quad
 \operatorname{vol}\bigl((C\setminus(C\ominus[-w,w]^m))\cap Q\bigr)\le2mw(2\sigma)^{m-1}.
\]

\emph{Tube containment.} Let $w>\gamma$ and let
$\operatorname{dist}_\infty(x,\partial C)\le\gamma$. Since $\partial C$ is
closed, there is $y\in\partial C$ with $\norm{x-y}_\infty\le\gamma$. As
$y\in C$, we get $x\in C+[-w,w]^m$. As $y$ is a boundary point, there is
an exterior point $u\notin C$ with $\norm{u-y}_\infty<w-\gamma$; then
$\norm{u-x}_\infty<w$, so $x+[-w,w]^m\not\subseteq C$ and
$x\notin C\ominus[-w,w]^m$. Hence the tube lies in the union of the two
shells. Dividing their total volume bound $4mw(2\sigma)^{m-1}$ by
$\operatorname{vol}(Q)=(2\sigma)^m$ gives probability at most $2mw/\sigma$.
Letting $w\downarrow\gamma$ proves \eqref{eq:cubetube}, including
$\gamma=0$. Only closedness and convexity of $C$ were used, so unbounded
and lower-dimensional sets are covered.

Finally, $\norm{\cdot}_\infty\le\norm{\cdot}_2$, so the Euclidean event
is contained in the infinity-norm event.
\end{proof}
The coefficient $2m$ is sharp for bounds of the form $c_m\gamma/\sigma$.
Let $C=b_0+[-a,a]^m$ with $0<a<\sigma$, and let $0<\gamma\le a$ with
$a+\gamma\le\sigma$. The tube is the cube of radius $a+\gamma$ minus the
open cube of radius $a-\gamma$, both centred at $b_0$ and contained in $Q$,
so its probability is $((a+\gamma)^m-(a-\gamma)^m)/\sigma^m$. Its
derivative at $\gamma=0$ is $(2m/\sigma)(a/\sigma)^{m-1}$, which tends to
$2m/\sigma$ as $a\uparrow\sigma$.

\subsection{Rational-grid exactness and its encoding}
To obtain rational right-hand sides of controlled encoding length, we
sample from a finite grid and couple the grid sample with a uniform point
of $Q$. For an integer $q\ge1$, let $h=2\sigma/q$ and sample each
coordinate of $G$ independently and uniformly from the $q$ cell centres
\begin{equation}\label{eq:centergrid}
 b_{0i}-\sigma+(j+1/2)h,\qquad j=0,\ldots,q-1.
\end{equation}
Let $U$ be uniform on $[-h/2,h/2]^m$ and independent of $G$. Then
$B=G+U$ is uniform on $Q$, and $\norm{B-G}_\infty\le h/2=\sigma/q$.
Since $\operatorname{dist}_\infty(\cdot,\partial C)$ is one-Lipschitz in
the infinity norm,
\[
 \{\operatorname{dist}_\infty(G,\partial C)\le\gamma\}\subseteq
 \{\operatorname{dist}_\infty(B,\partial C)\le\gamma+\sigma/q\},
\]
and Lemma~\ref{lem:cubetube} gives, for every closed convex $C$ and
$\gamma\ge0$,
\begin{equation}\label{eq:gridtube}
 \Pr\{\operatorname{dist}_\infty(G,\partial C)\le\gamma\}
       \le\min\{1,2m(\gamma/\sigma+1/q)\}.
\end{equation}
The rounding term $2m/q$ is necessary. Convex boundaries have Lebesgue
measure zero, but grid points can lie on them with positive probability:
for $m=1$, an interval whose two endpoints are distinct grid points has
boundary probability $2/q$ at $\gamma=0$, exactly the correction in
\eqref{eq:gridtube}.

\begin{theorem}\label{thm:gridsmoothing}
Assume the setting of Section~\ref{sec:perturbation-setting}. Let
$b_0\in\Q^m$ and let $\sigma,\epsilon,\underline f,\overline f$ be
rational with $\sigma>0$ and $0<\epsilon<1$. Let $G$ be sampled from the
grid \eqref{eq:centergrid} with
\begin{equation}\label{eq:gridsmoothing}
 q\ge\left\lceil\frac{4mK}{\epsilon}\right\rceil,\qquad
 \gamma=\frac{\sigma\epsilon}{4mK},\qquad
 \bar\rho=\frac M\gamma=\frac{4mKM}{\sigma\epsilon}.
\end{equation}
Then, with probability at least $1-\epsilon$, either the sampled problem
is infeasible or $p_{\bar\rho}(G)=v(G)$. On the same event, if the
sampled problem is feasible, every coefficient $\rho>\bar\rho$ is
minimizer-exact at multiplier zero.
\end{theorem}
\begin{proof}
By \eqref{eq:gridtube} and a union bound over the at most $K$ images,
\[
 \Pr\{\operatorname{dist}_\infty(G,\partial C_z)\le\gamma
      \text{ for some }z\in\mathcal Z\}
 \le2mK\Bigl(\frac\gamma\sigma+\frac1q\Bigr)
 \le2mK\Bigl(\frac\epsilon{4mK}+\frac\epsilon{4mK}\Bigr)=\epsilon.
\]
On the complementary event every boundary distance exceeds $\gamma$, and
Lemma~\ref{lem:rhsmargin} applies whenever $G$ is feasible.
\end{proof}
In words: with probability at least $1-\epsilon$, the grid sample has
boundary distance greater than $\gamma=\sigma\epsilon/(4mK)$ from every
residual image; then either no slice is feasible, or
$\rho_*[G]\le\bar\rho=M/\gamma$.

\begin{remark}[Encoding]\label{rem:gridencoding}
Choose $q$ as the smallest power of two with
$q\ge\lceil4mK/\epsilon\rceil$. Then
$\log_2q=O(1+\log m+\log K+\log(1/\epsilon))$, and exact uniform sampling
uses $m\log_2q$ independent random bits. The grid points and $\bar\rho$
have encoding length polynomial in the supplied rational data
$b_0,\sigma,\epsilon,\underline f,\overline f$, in $m$ and in $\log K$; an
integer coefficient strictly larger than $\bar\rho$, such as
$\lfloor\bar\rho\rfloor+1$, has the same bound, although $\bar\rho$ is
numerically of order $K$. Neither computing $\bar\rho$ nor sampling
requires enumerating the slices. The construction does not check the
structural assumptions of Section~\ref{sec:perturbation-setting} and
does not solve the sampled problem.
\end{remark}

\begin{remark}[Endpoint grid]\label{rem:endpointgrid}
Theorem~\ref{thm:gridsmoothing} also holds for the grid with $q+1$ points
$b_{0i}-\sigma+jh$, $0\le j\le q$, per coordinate. Its jitter cells tile
the cube of radius $\sigma+h/2$, so Lemma~\ref{lem:cubetube} for that cube
gives
\[
 \Pr\{\operatorname{dist}_\infty(G,\partial C)\le\gamma\}
 \le\min\left\{1,\frac{2m(\gamma+h/2)}{\sigma+h/2}\right\}
 \le\min\{1,2m(\gamma/\sigma+1/q)\}.
\]
With $q+1$ points per coordinate, the random-bit count of
Remark~\ref{rem:gridencoding} no longer applies. For either grid the
Euclidean boundary event obeys the same estimate, but it pairs only with
a Euclidean penalty: a Euclidean margin greater than $\gamma$ does not
give an infinity-norm margin greater than $\gamma$.
\end{remark}

\subsection{Feasibility, local stability and the original instance}
\label{sec:perturbation-stability}
\paragraph{The infeasibility alternative.}
Theorem~\ref{thm:gridsmoothing} cannot guarantee feasibility. For
example, if no residual image meets $Q$, every sample is infeasible. If
some slice satisfies $Q\subseteq C_{z_0}$, every grid point is feasible,
and the theorem gives exactness at multiplier zero with probability at
least $1-\epsilon$ unconditionally. More generally, if the grid sample is
feasible with probability at least $p_0>0$, then the conditional failure
probability given feasibility is at most $\epsilon/p_0$.

This bound concerns the grid actually used; refining the grid can lower
its feasibility probability. For a recipe valid for every grid size,
suppose that $B$ uniform on $Q$ is feasible with probability at least a
known rational $p_0>0$. Under the jitter coupling, if $B\in C_z$ and $G\notin C_z$,
then the segment from $G$ to $B$ crosses $\partial C_z$, so
$\operatorname{dist}_\infty(B,\partial C_z)\le\sigma/q$.
Lemma~\ref{lem:cubetube} and a union bound give
\[
 \Pr\{G\text{ feasible}\}\ge p_0-\frac{2mK}q .
\]
Running the construction with $\epsilon p_0/2$ in place of $\epsilon$
(so $q\ge\lceil8mK/(\epsilon p_0)\rceil$ and $\bar\rho=8mKM/(\sigma\epsilon p_0)$)
gives $\Pr\{G\text{ feasible}\}\ge p_0(1-\epsilon/4)$, while the failure
probability is at most $\epsilon p_0/2$. Hence the conditional failure
probability given feasibility is at most
$\epsilon/(2-\epsilon/2)<\epsilon$.

\paragraph{Local stability.}
Let $b$ be feasible with
$\operatorname{dist}_\infty(b,\partial C_z)>\gamma$ for every $z$. By the
first step of the proof of Lemma~\ref{lem:rhsmargin}, every residual
image $C_z$ that contains $b$ contains the whole cube
$b+[-\gamma,\gamma]^m$, while every other image has distance greater than
$\gamma$ from $b$. Hence the set of slices (integer assignments) that are
feasible for the right-hand side is the same at every point of the open cube
of radius $\gamma$ about $b$. We use the smaller cube
$b+[-\gamma/2,\gamma/2]^m$ because, by the one-Lipschitz property of the
distance, every point $b'$ of it keeps margin greater than $\gamma/2$.
Applying the repair inequality of Lemma~\ref{lem:rhsmargin} at $b'$ with
margin $\gamma/2$ to an optimal point for $b''$, and then with the roles
exchanged, gives
\[
 |v_z(b')-v_z(b'')|\le\frac{2M}\gamma\norm{b'-b''}_\infty
\]
for every slice $z$ whose image contains $b$ and all $b',b''$ in the
cube. A minimum over
the same finite collection of slice values keeps this Lipschitz bound, so
$v$ is $2M/\gamma$-Lipschitz in the infinity norm on the cube. By
Lemma~\ref{lem:rhsmargin} with margin $\gamma/2$, the coefficient
$2M/\gamma$ is exact at multiplier zero at every point of the cube, and
every larger coefficient is minimizer-exact at multiplier zero there.
Lipschitz value bounds across discrete choices under uniform Slater
conditions also appear in mixed-integer optimal control
\citep[Theorem~3]{GugatHante2017}.

\paragraph{The original instance.}
The local conclusion does not recover the optimum at $b_0$. The chain of
Section~\ref{sec:lower} fits the present setting with $m=1$, two slices
and $M=1$: the slice with binary variable zero has residual image
$[-1,1]$ and cost $0$, and the improving slice (binary variable one) has
residual image $[\delta_d,1]$ and cost $-1$. The right-hand side $0$ lies
within $\delta_d$ of the boundary of $[\delta_d,1]$, so
Lemma~\ref{lem:rhsmargin} gives only coefficients above $1/\delta_d$,
the zero-multiplier threshold of Section~\ref{sec:lower}; this is why the
coefficient explodes. Moving the right-hand side to $\delta_d$ makes the
improving assignment feasible, and the optimum becomes $-1$ instead of
$0$. Short penalties after perturbation can therefore come with a
different optimal integer assignment.

\subsection{Numerical size: unavoidable tails and slice dependence}
\label{sec:numerical-size}
The coefficient $\bar\rho=4mKM/(\sigma\epsilon)$ of
Theorem~\ref{thm:gridsmoothing} has short encoding but can be
numerically large. The following examples show which parts of this size
are unavoidable:
\begin{itemize}
 \item[(a)] the order $K/\epsilon$ is necessary: some instances satisfy
 $\Pr\{\rho_*[B]>T\}\ge(K-1)/T$ for $T\ge K$
 (Example~\ref{ex:tail});
 \item[(b)] a lower bound of order $K$ can hold almost surely, not just
 with small probability (Example~\ref{ex:slicecount});
 \item[(c)] two slices can give $\mathbb E\rho_*[B]=+\infty$ while
 $\mathbb E\log(1+\rho_*[B])$ is finite (Example~\ref{ex:infinitemean});
 \item[(d)] an exceptional grid point can have infinite threshold, so
 the high-probability guarantee cannot be turned into a bound on
 expectations (Example~\ref{ex:atom}).
\end{itemize}
In all examples $m=1$, so the norm is the absolute value, and
$\underline f=-1$, $\overline f=0$, $M=1$. All thresholds are optimized
thresholds $\rho_*[b]$ unless we say otherwise; since exactness at
multiplier zero implies value-exactness, lower bounds on $\rho_*[b]$ also
bound the zero-multiplier threshold. The examples list their slices
explicitly, as the slice model allows; they are not derived from a
succinct mixed-integer encoding.

The lower bounds use one balanced mixture
(Proposition~\ref{prop:mixture}). Suppose a slice with cost $-1$ has a
singleton residual image at distance $t>0$ from $b$, and a cost-zero
slice contains a point whose residual $r_z(x)-b$ has the opposite sign
and magnitude $s>0$. Put weight
$s/(s+t)$ on the singleton point and weight $t/(s+t)$ on the cost-zero
point. The residuals cancel, and the average penalized value is
\begin{equation}\label{eq:tailmixture}
 \frac{s}{s+t}(-1+2\rho t)
\end{equation}
for every multiplier. If $v(b)=0$, this mixture shows
$\rho_*[b]\ge1/(2t)$.

\begin{example}[The order $K/\epsilon$ is necessary]\label{ex:tail}
Let an anchor slice have residual image $[0,1]$ and cost zero, and let
$K-1$ singleton slices, $K\ge2$, have images $\{j/K\}$,
$1\le j\le K-1$, and cost $-1$. At every right-hand side
$b\in(0,1)$ that is not one of the points $j/K$, the primal value is
zero, and \eqref{eq:tailmixture} with the nearest singleton at distance
$t$ gives $\rho_*[b]\ge1/(2t)$. Let $B$ be uniform on $[0,1]$. For
$T\ge K$, the intervals $|B-j/K|<1/(2T)$ are disjoint, lie in $(0,1)$ and
have total length $(K-1)/T$. Consequently
\begin{equation}\label{eq:penaltytail}
 \Pr\{\rho_*[B]>T\}\ge\frac{K-1}T\qquad(T\ge K).
\end{equation}
Let $0<\epsilon<(K-1)/K$. For $K\le T<(K-1)/\epsilon$,
\eqref{eq:penaltytail} gives $\Pr\{\rho_*[B]>T\}>\epsilon$, and for
$T<K$ the same holds by monotonicity. Hence the $(1-\epsilon)$-quantile
of $\rho_*[B]$ is at least $(K-1)/\epsilon$. Here $b_0=1/2$,
$\sigma=1/2$ and $M=1$, so Theorem~\ref{thm:gridsmoothing} gives the
coefficient $8K/\epsilon$; the order $K/\epsilon$ cannot be improved, even
though every sample is feasible and the multiplier is optimized.

The grid sample of Theorem~\ref{thm:gridsmoothing} behaves the same way.
On $[0,1]$ the cell centres are $(i+1/2)/q$; let $K$ divide $q$, so that
no centre is a singleton image. If $q\ge2K$, each singleton image $j/K$
has two adjacent centres at distance $1/(2q)$, and these $2(K-1)$ centres
are distinct; hence $\Pr\{\rho_*[G]\ge q\}\ge2(K-1)/q$. More generally,
for $K\le T\le q$, each singleton image has at least $q/T-1$ centres at
distance less than $1/(2T)$, these sets are disjoint for different $j$,
and at each such centre $\rho_*[G]>T$. Hence
$\Pr\{\rho_*[G]>T\}\ge(K-1)(1/T-1/q)$. If $\epsilon\le(K-1)/(2K)$ and
$q\ge4K/\epsilon$, then $T=(K-1)/(2\epsilon)$ satisfies $K\le T\le q$, and
$\Pr\{\rho_*[G]>T\}\ge2\epsilon-(K-1)/q>\epsilon$. So the
$(1-\epsilon)$-quantile of $\rho_*[G]$ is at least $(K-1)/(2\epsilon)$,
and the order $K/\epsilon$ is necessary for the grid result as well.
\end{example}

\begin{example}[An almost-sure slice-count obstruction]\label{ex:slicecount}
Keep the anchor slice with image $[0,1]$ and cost zero, and use $J+1$
singleton slices, $J\ge1$, with images $\{j/J\}$, $j=0,\ldots,J$, and
cost $-1$. There are $K=J+2$ slices. Let $b\in(0,1)$ lie strictly between
the neighbouring singleton images $j/J$ and $(j+1)/J$, and let
$g_1=b-j/J>0$ and $g_2=(j+1)/J-b>0$ be its distances to them, so
$g_1+g_2=1/J$. The primal value is zero. Put weight $g_2/(g_1+g_2)$ on
the lower singleton and $g_1/(g_1+g_2)$ on the upper one. The residuals
cancel, the cost is $-1$, and the mean absolute residual is
\[
 \frac{2g_1g_2}{g_1+g_2}=2Jg_1g_2\le\frac1{2J}.
\]
Hence $\rho_*[b]\ge2J=2(K-2)$ for every $b\in(0,1)$ other than the
singleton images, that is, for almost every $b\in[0,1]$.
\end{example}

\begin{example}[Infinite mean with two slices]\label{ex:infinitemean}
Take an anchor slice with image $[-1,1]$ and cost zero and a singleton
slice with image $\{0\}$ and cost $-1$. For $b\in(-1,1)\setminus\{0\}$,
the primal value is zero, and \eqref{eq:tailmixture} with $t=|b|$ gives
$\rho_*[b]\ge1/(2|b|)$. Conversely, at $\rho=1/(2|b|)$ and multiplier
$\lambda=-\rho\operatorname{sign}(b)$, the singleton's penalized value is
$-1+\lambda(-b)+\rho|b|=-1+2\rho|b|=0$, and the anchor's penalized value
$\lambda e+\rho|e|$ at residual $e$ is nonnegative. Hence
\begin{equation}\label{eq:infinitemean}
 \rho_*[b]=\frac1{2|b|}.
\end{equation}
The zero-multiplier threshold is instead $1/|b|$. For $B$ uniform on
$[-1,1]$, we get $\mathbb E\rho_*[B]=+\infty$, while
$\mathbb E\log(1+\rho_*[B])<\infty$ because $\log(1+1/(2|b|))$ is
integrable near zero.
\end{example}

\begin{example}[An exceptional grid point with infinite threshold]
\label{ex:atom}
Take a convex quadratic slice
$\{(\alpha,\beta):0\le\alpha\le1,\ \alpha^2\le\beta\le1\}$ with cost
$-\alpha$ and residual $\beta$, so its image is $[0,1]$, and a singleton
slice with image $\{-1\}$ and cost zero. At right-hand side $0$ the
constraint $\beta=0$ forces $\alpha=0$, so the primal value is zero. For
$0<u\le1$, mix the point $(\alpha,\beta)=(u,u^2)$ with weight
$1/(1+u^2)$ and the singleton with weight $u^2/(1+u^2)$. The residuals
cancel, the cost is $-u/(1+u^2)$, and the mean absolute residual is
$2u^2/(1+u^2)$. By Corollary~\ref{cor:ratio}, $\rho_*[0]$ is at least the
ratio $1/(2u)$ for every such $u$, so $\rho_*[0]=+\infty$: no finite
penalty is value-exact. A rational grid containing $0$, for instance the
grid \eqref{eq:centergrid} with $b_0=0$ and odd $q$, gives this
right-hand side positive probability. Hence $\mathbb E\rho_*[G]=+\infty$
(with any nonnegative value assigned to infeasible samples), and no
almost-sure bound on the threshold or its encoding holds. The margin
condition of Lemma~\ref{lem:rhsmargin} excludes this point, since $0$
lies on the boundary of $[0,1]$. The refined Slater condition also fails
there: on $\beta=0$ the inequality $\alpha^2\le\beta$ cannot hold strictly.
\end{example}

\subsection{A Gaussian companion in the Euclidean norm}\label{sec:gaussian}
Lemma~\ref{lem:rhsmargin} and its proof hold for any norm on $\R^m$,
provided the same norm defines both the margin and the penalty: replace
the cube $b+[-\gamma,\gamma]^m$ by the ball of radius $\gamma$. For
Gaussian noise we use the Euclidean norm for both. There is no grid
here, so this result has no encoding statement.

\begin{proposition}\label{prop:gaussian}
Assume the setting of Section~\ref{sec:perturbation-setting}, but with
the Euclidean norm in the penalty of \eqref{eq:rhsmodel}. Let
$B=b_0+\sigma\xi$ with $\xi\sim\mathcal N(0,I_m)$, $\sigma>0$ and
$0<\epsilon<1$. Then for every $\gamma\ge0$,
\begin{equation}\label{eq:gaussiantube}
 \Pr\{\min_{z\in\mathcal Z}\operatorname{dist}_2(B,\partial C_z)\le\gamma\}
       \le\frac{8Km^{1/4}\gamma}\sigma .
\end{equation}
Consequently, with probability at least $1-\epsilon$, either the sampled
problem is infeasible or the coefficient
\begin{equation}\label{eq:gaussianpenalty}
 \frac{8MKm^{1/4}}{\sigma\epsilon}
\end{equation}
is exact at multiplier zero for the Euclidean penalty; every larger
coefficient is then minimizer-exact at multiplier zero.
\end{proposition}
\begin{proof}
First let $C$ be a nonempty compact convex set with nonempty interior.
\citet[Lemma~2.3.4]{DunaganSpielmanTeng2011} bound the probabilities of
the outer tube $\{\operatorname{dist}_2(B,\partial C)\le\gamma\}\setminus C$
and of the inner tube $\{\operatorname{dist}_2(B,\partial C)\le\gamma\}\cap C$
each by $4m^{1/4}\gamma/\sigma$. Hence
$\Pr\{\operatorname{dist}_2(B,\partial C)\le\gamma\}\le8m^{1/4}\gamma/\sigma$.

Now let $C$ be nonempty, compact and convex with empty interior, so
$\partial C=C$. For $\eta>0$, the parallel body
$C_\eta=C+\eta\mathbb B_2$ is compact and convex with nonempty interior.
We claim $\operatorname{dist}_2(x,\partial C_\eta)\le|\operatorname{dist}_2(x,C)-\eta|$.
Let $p$ be the projection of $x$ onto $C$, and let $u$ be a unit vector
in the normal cone of $C$ at $p$: take $u=(x-p)/\norm{x-p}_2$ if
$x\notin C$, and a unit normal to the affine hull of $C$ if $x\in C$
(it exists because $C$ has empty interior). Then $p+su$ has projection
$p$ and distance $s$ from $C$ for every $s\ge0$. So $p+\eta u\in C_\eta$,
while $p+su\notin C_\eta$ for $s>\eta$; thus $p+\eta u\in\partial C_\eta$.
Since $x=p+\operatorname{dist}_2(x,C)\,u$, the claim follows. Hence
$\operatorname{dist}_2(x,C)\le\gamma$ implies
$\operatorname{dist}_2(x,\partial C_\eta)\le\gamma+\eta$, and the
first case gives
\[
 \Pr\{\operatorname{dist}_2(B,\partial C)\le\gamma\}
 \le\Pr\{\operatorname{dist}_2(B,\partial C_\eta)\le\gamma+\eta\}
 \le\frac{8m^{1/4}(\gamma+\eta)}\sigma .
\]
Letting $\eta\downarrow0$ gives the same bound as for convex bodies.

A union bound over the at most $K$ images proves \eqref{eq:gaussiantube}.
With $\gamma=\sigma\epsilon/(8Km^{1/4})$, the right-hand side equals
$\epsilon$. On the complementary event, every Euclidean boundary distance
exceeds $\gamma$, and the Euclidean version of Lemma~\ref{lem:rhsmargin}
shows that $M/\gamma$, which equals \eqref{eq:gaussianpenalty}, is exact
at multiplier zero whenever $B$ is feasible, and that every larger
coefficient is minimizer-exact at multiplier zero.
\end{proof}

\section{Discussion}\label{sec:discussion}
The main conclusion is a limit on extending polynomial penalty encoding
from linear to convex quadratic native constraints. Finite exactness and
uniform strict feasibility on feasible integer slices do not suffice:
an improving infeasible slice can lie doubly exponentially close to the
linking equality. The upper bounds identify two ways to recover
polynomial encoding, by fixing continuous dimension or the number of
nonlinear native inequalities under the stated regularity assumptions.
These results refine the quantitative scope of existing exactness
theorems; they do not replace their existence theory.

The calibration result then shows why a safe sufficient coefficient and
a nearly smallest one are different computational targets, even when
the original optimizer is supplied. The perturbation result provides a
probabilistic alternative with explicit sampling and feasibility
qualifications. The encoding upper bounds and the perturbation guarantee
allow large numerical coefficients and do not imply fast solution of the
augmented subproblem. Changing the equality, using succinct exponents or
choosing a fractional-power augmentation changes the representation
question. The contributions concern
ordinary binary coefficients for the prescribed norm and residual.

Supplementary scripts check the exact finite envelopes, determinant
identities, degenerate regularization, calibration dual formulas and thresholds, and selected
grid and tail identities with rational or symbolic arithmetic. These
checks support the algebra and indexing; the proofs establish the general
claims. No solver-performance claim is made.

\bibliographystyle{plainnat}
\bibliography{references}
\appendix
\section{An exact dual value without an attaining multiplier}\label{app:attainment}
Example~13 of the December 15, 2025 manuscript of
\citet{LefebvreSchmidt2025} is stated to admit no finite exact penalty
representation. Under its supremum-based Definition~1, however, the
example has an exact dual value at every nonnegative penalty.
In renamed variables its native set is
\[
 X=\{(u,s,z)\in[-1,1]^2\times\{0,1\}:u+2z\le1,\ u^2\le s\},
 \qquad f=u-z/2,\qquad r=s.
\]
The equality $s=0$ forces $u=z=0$, so $v=0$. All native residuals are
nonnegative. Proposition~\ref{prop:mixture} therefore gives $D_\rho=0$
for every $\rho\ge0$.

This can also be checked directly. Put $t=\lambda+\rho$ and assume
$t\ge1/2$. The branch $z=1$ contains only $(-1,1,1)$ and has value
$t-3/2$. On $z=0$, minimization first sets $s=u^2$ and then
$u=-1/(2t)$, giving
\[
 L_\rho(\lambda)=\min\{t-3/2,-1/(4t)\}.
\]
This tends to zero as $\lambda\to+\infty$. For any finite $t$, however,
a sufficiently small $\eta>0$ gives the native point
$(-\eta,\eta^2,0)$ with value $-\eta+t\eta^2<0$.
Thus no finite multiplier attains the exact dual value. The source's reduction to multiplier zero invokes its Lemma~3,
which transfers attained exactness between fixed multipliers and does not
apply to this nonattained supremum. Thus the example demonstrates
nonattainment, not failure of dual-value exactness. This correction does
not affect the source's Theorem~14 under slice Slater conditions, nor the
attained lower-bound formulas in Section~\ref{sec:lower}.

For genuine failure of finite dual-value exactness without refined Slater,
the two-sided exceptional-atom construction in Section~\ref{sec:perturbation}
has an infinite optimized threshold at right-hand side zero.

\section{A symmetric two-binary variant}\label{app:symmetric}
In the instance of Theorem~\ref{thm:lower}, the improving residuals near
zero are all positive, so a positive multiplier halves the required
coefficient. This appendix gives a variant in which residuals of both
signs occur at cost $-1$. Any nonzero multiplier then hurts one of the two
sides, the optimal multiplier is zero, and the optimized and
zero-multiplier thresholds coincide at $2^{2^d}$.

Keep the chain $A_d$ and $\delta_d=2^{-2^d}$ of~\eqref{eq:chain}, the
objective $-q$ and the relaxed equality $y=0$. Add a second binary
variable $b\in\{0,1\}$ and replace the switching inequality
in~\eqref{eq:onebinary} by the pair
\begin{equation}\label{eq:symmetric}
 y\ge a_d-(1-q)-2(1-b),\qquad
 y\le-a_d+(1-q)+2b.
\end{equation}
The native set is thus
$\{(a,y,q,b):a\in A_d,\ -1\le y\le1,\ q,b\in\{0,1\},\ \eqref{eq:symmetric}\}$.
At $q=1$, the choice $b=1$ forces $y\ge a_d$ and the choice $b=0$ forces
$y\le-a_d$, so the improving slices lie on both sides of $y=0$. We write
$L^{\rm sym}_{d,\rho}(\lambda)$ and $D^{\rm sym}_{d,\rho}$ for the
penalized and dual values with the absolute-value penalty.

\begin{proposition}\label{prop:symmetric}
Let $d\ge1$. The variant has optimal value $0$. For every $\rho\ge0$ and
$\lambda\in\R$, with $c=\rho-|\lambda|$,
\begin{equation}\label{eq:symdual}
 L^{\rm sym}_{d,\rho}(\lambda)=\min\{0,\,-1+c\delta_d,\,-1+c\},\qquad
 D^{\rm sym}_{d,\rho}=\min\{0,\,-1+\rho\delta_d\},
\end{equation}
and $\lambda=0$ attains the dual value. Consequently:
\begin{enumerate}
\item[(a)] the least penalty and the zero-multiplier threshold both equal
$\delta_d^{-1}=2^{2^d}$, which has $2^d+1$ bits, and at
$\rho=\delta_d^{-1}$ the multiplier $\lambda=0$ is the only one at which
$\rho$ is exact;
\item[(b)] for $0\le\epsilon<1$, the gap $0-D^{\rm sym}_{d,\rho}$ is at
most $\epsilon$ exactly when $\rho\ge(1-\epsilon)\delta_d^{-1}$;
\item[(c)] every $\rho>\delta_d^{-1}$ is minimizer-exact at multiplier
zero;
\item[(d)] every integer slice $(q,b)$ contains a point at which every
native inequality in the continuous variables holds with slack at least
$1/8$; for $q=0$ this point also satisfies $y=0$.
\end{enumerate}
\end{proposition}
\begin{proof}
As for Theorem~\ref{thm:lower}, we first compute the projection onto
$(q,y)$, on which the objective and the residual depend, and then minimize
over it.

\emph{Projection.} Fix $a\in A_d$. Using $-1\le y\le1$ and
$\delta_d\le a_d\le1$, the four slices give the following $y$-intervals:
$[a_d,1]$ for $(q,b)=(1,1)$, since the upper bound $2-a_d$ is at least
one; $[-1,-a_d]$ for $(1,0)$, since the lower bound $a_d-2$ is at most
$-1$; $[a_d-1,1]$ for $(0,1)$; and $[-1,1-a_d]$ for $(0,0)$. At $q=1$
every native point therefore has $|y|\ge a_d\ge\delta_d$. At $a=\bar a$
the two intervals at $q=1$ are $[\delta_d,1]$ and $[-1,-\delta_d]$, and
the two intervals at $q=0$ are $[\delta_d-1,1]$ and $[-1,1-\delta_d]$,
which cover $[-1,1]$. Hence the projection is
\[
 (\{0\}\times[-1,1])\ \cup\
 (\{1\}\times([-1,-\delta_d]\cup[\delta_d,1])).
\]
Since $\delta_d>0$, the relaxed equality $y=0$ forces $q=0$, and the
optimal value is $0$.

\emph{Penalized value.} For a point with residual $y$, the term
$\lambda y+\rho|y|$ is at least $c|y|$, with equality for one sign of $y$.
Both signs occur on each part of the projection, and each part is
symmetric in $y$, so the minimum of $\lambda y+\rho|y|$ over a part equals
the minimum of $cs$ over the attainable magnitudes $s=|y|$. At $q=0$ these
are $s\in[0,1]$, giving $\min\{0,c\}$. At $q=1$ they are
$s\in[\delta_d,1]$, giving $-1+\min\{c\delta_d,c\}$. The term $c$ is
larger than $-1+c$, so it can be dropped, which proves the first formula
in~\eqref{eq:symdual}. This expression is nondecreasing in $c$, and
$c\le\rho$ with equality only at $\lambda=0$. Hence the dual value is
attained at $\lambda=0$ and equals $\min\{0,-1+\rho\delta_d,-1+\rho\}$.
Since $\delta_d\le1$, the last term can be dropped, which proves the
second formula.

\emph{Parts (a)--(c).} By~\eqref{eq:symdual}, $D^{\rm sym}_{d,\rho}=0$
exactly when $\rho\ge\delta_d^{-1}$, and the same holds for
$L^{\rm sym}_{d,\rho}(0)$. At $\rho=\delta_d^{-1}$ and $\lambda\ne0$ we
have $c<\delta_d^{-1}$, so $-1+c\delta_d<0$. The integer $2^{2^d}$ has
$2^d+1$ bits. For (b), the gap is $\max\{0,1-\rho\delta_d\}$, which is at
most $\epsilon\ge0$ exactly when $\rho\ge(1-\epsilon)\delta_d^{-1}$. For
(c), take $\lambda=0$ and $\rho>\delta_d^{-1}$. A native point with $q=1$
has penalized value $-1+\rho|y|\ge-1+\rho\delta_d>0$, and a native point
with $q=0$ and $y\ne0$ has penalized value $\rho|y|>0$. So the minimizers
are exactly the feasible points.

\emph{Part (d).} Set $a_i=3/8$ for all $i$. As in the proof of
Theorem~\ref{thm:lower}, the chain and its bounds have slack at least
$1/8$. On the slices with $q=0$ take $y=0$: the inequalities
in~\eqref{eq:symmetric} have slacks $21/8$ and $5/8$ for $b=0$, and $5/8$
and $21/8$ for $b=1$. On the slices with $q=1$ take $y=(2b-1)/2$: for
$b=1$ the slacks are $1/8$ and $9/8$, and for $b=0$ they are $9/8$ and
$1/8$. The bounds on $y$ have slack at least $1/2$. The smallest slack is
$1/8$.
\end{proof}

At $\rho=\delta_d^{-1}$ and $\lambda=0$, the points $(q,y)=(1,\pm\delta_d)$
have penalized value $0$, so this coefficient is exact but not
minimizer-exact. The variant has two binary variables, $d+1$ continuous
variables, a linear objective and the same chain of $d-1$ convex
quadratic inequalities as the instance of Theorem~\ref{thm:lower}. The
fractional-power augmentation of Section~\ref{sec:lower-scope} behaves
as before on this variant, because $|y|\ge\delta_d$ on its slices with
$q=1$: the coefficient two is exact, and every larger coefficient is
minimizer-exact at multiplier zero.

\end{document}
